\documentclass[11pt]{article}

\usepackage[margin=1in]{geometry}
\usepackage{graphicx}
\usepackage{subcaption}
\usepackage{booktabs}
\usepackage{longtable}
\usepackage{array}
\usepackage[hidelinks]{hyperref}
\usepackage{float}                % enables [H] placement
\usepackage[section]{placeins}    % prevents floats crossing sections

\title{WESAD: Data Transformation and Exploratory Data Analysis}
\author{Anishwar Chakraborty}
\date{\today}

% --- helpers ---
\newcommand{\FigIfExists}[4]{%
  \begin{figure}[H]
    \centering
    \IfFileExists{#1}{\includegraphics[width=#2\linewidth]{\detokenize{#1}}}{\fbox{\parbox{0.95\linewidth}{Missing figure: \texttt{#1}\\\newline Run the corresponding EDA script to generate it.}}}
    \caption{#3}
    \label{#4}
  \end{figure}
}

\newcommand{\SubFigIfExists}[5]{%
  \begin{subfigure}{#2\linewidth}
    \centering
    \IfFileExists{#1}{\includegraphics[width=\linewidth]{\detokenize{#1}}}{\fbox{\parbox{0.95\linewidth}{Missing figure: \texttt{#1}}}}
    \caption{#3}
    \label{#4}
  \end{subfigure}
}

\begin{document}
\maketitle

\section{WESAD overview (what it is and what it contains)}
WESAD (Wearable Stress and Affect Detection) is a multimodal wearable-sensing dataset collected in a controlled lab protocol to study how
physiological and movement signals relate to stress and affect.
Each subject experiences protocol segments such as \textbf{baseline} (rest), \textbf{stress} (stress induction), and \textbf{amusement} (positive affect).

\subsection{What data types are present?}
WESAD is not an image or text dataset. It contains:
\begin{itemize}
  \item \textbf{Time-series sensor signals} from wrist and chest devices (e.g., EDA, BVP/PPG, ACC, temperature, respiration).
  \item \textbf{A synchronized label stream} indicating which protocol segment is active over time.
  \item \textbf{Subject metadata} (demographics and questionnaire-style yes/no responses) stored in per-subject readme files.
\end{itemize}

\subsection{Why it is useful for stress modeling}
WESAD supports stress detection research because it enables:
\begin{itemize}
  \item \textbf{Multimodal learning}: combining physiological and movement signals.
  \item \textbf{Window-based supervised learning}: converting continuous recordings into labeled examples.
  \item \textbf{Cross-subject generalization}: evaluating whether a model trained on some subjects works on new subjects.
\end{itemize}

\section{Data transformation (implemented in \texttt{data\_wrangling.py})}
Raw WESAD recordings are multi-rate time series (different sensors have different sampling frequencies).
Most EDA and most classifiers are easier to run on a \textbf{fixed-length window} dataset where each row is one labeled example.
The transformation pipeline converts raw signals into a window-level feature table.

\subsection{What the transformation does (conceptually)}
\begin{enumerate}
  \item \textbf{Load raw wrist/chest signals and labels.}
  \item \textbf{Denoise signals} (e.g., smoothing / low-pass filtering) to reduce high-frequency noise and improve robustness.
  \item \textbf{Time-align modalities.} Because sensors are sampled at different rates, signals are aligned onto a consistent timeline.
        Labels are propagated so every aligned time point (and later every window) has an associated class.
  \item \textbf{EDA decomposition (cvxEDA).} EDA can be decomposed into:
        \textbf{tonic} (slow baseline drift) and \textbf{phasic} (fast skin conductance responses), which can be more informative for stress.
  \item \textbf{Split into protocol segments} using the synchronized label stream.
  \item \textbf{Windowing.} The continuous signal is partitioned into fixed windows (30 seconds), creating labeled examples.
  \item \textbf{Feature extraction per window.} For each sensor channel, compute summary statistics (mean/std/min/max) and a few derived features.
\end{enumerate}

\subsection{Why summary statistics are used as features}
Window features compress high-frequency signals into stable vectors that still capture the window's overall level, variability, and extremes.

\subsubsection*{Mean, median, and mode (what they are)}
\begin{itemize}
  \item \textbf{Mean}: arithmetic average; captures typical level but is sensitive to spikes/outliers.
  \item \textbf{Median}: middle value; robust to spikes (often useful when artifacts are common).
  \item \textbf{Mode}: most frequent value; usually more useful for discrete variables (e.g., labels) than for continuous physiology.
\end{itemize}

\subsubsection*{Other window statistics used here}
\begin{itemize}
  \item \textbf{Standard deviation (std)}: within-window variability (often rises with motion artifacts in BVP).
  \item \textbf{Min/max}: extremes (physiological peaks or artifacts).
  \item \textbf{Slope}: trend across a window (e.g., temperature drift).
  \item \textbf{Peak frequency}: dominant periodic structure (useful for BVP).
  \item \textbf{ACC magnitude}: compact movement intensity (also highlights confounding).
\end{itemize}

\subsection{Sensor modalities and derived features (plain-language)}
\begin{itemize}
  \item \textbf{EDA}: skin conductance; reflects sympathetic arousal, but drifts slowly and depends on skin contact.
  \item \textbf{EDA phasic}: fast-changing EDA component (SCR-like responses); often more directly linked to arousal.
  \item \textbf{EDA tonic}: slow EDA baseline level (SCL); captures longer-term shifts and drift.
  \item \textbf{EDA SMNA / sparse driver}: decomposition component representing estimated ``event-like'' driver activity.
  \item \textbf{BVP (PPG)}: optical pulse signal; carries heart-beat periodicity but is very motion-sensitive on the wrist.
  \item \textbf{BVP peak frequency}: dominant BVP frequency in a window; proxy for periodic structure/heart-rate band behavior (can be distorted by artifacts).
  \item \textbf{TEMP}: wrist skin temperature; slow-varying, often informative via trends rather than rapid changes.
  \item \textbf{TEMP slope}: temperature trend within a window (warming/cooling).
  \item \textbf{ACC}: 3-axis accelerometer; movement intensity/posture changes; important both as predictor and artifact indicator.
  \item \textbf{net\_acc}: acceleration magnitude (roughly $\sqrt{x^2+y^2+z^2}$) summarized over a window; compact movement feature.
  \item \textbf{Resp}: respiration signal; stress can affect breathing rate/variability and complements EDA/BVP.
\end{itemize}

\section{Metadata parsing and merging (implemented in \texttt{readme\_parser.py})}
Subjects differ in demographic and contextual factors (age, gender, smoking, etc.). These fields are extracted and encoded (e.g., one-hot indicators)
to support interpretation and confound analysis.

\section{Final dataset: \texttt{data/m14\_merged.csv}}
\subsection{What one row represents}
Each row corresponds to one \textbf{30-second window} from one subject, represented by engineered window features plus a class label and metadata.

\subsection{Label meanings}
The labels used in \texttt{m14\_merged.csv} are:
\begin{center}
\begin{tabular}{cl}
\toprule
\textbf{Label value} & \textbf{Meaning} \\
\midrule
0 & amusement \\
1 & baseline \\
2 & stress \\
\bottomrule
\end{tabular}
\end{center}

\subsection{3--4 sample rows (subset of columns)}
\begin{table}[H]
  \centering
  \small
  \begin{tabular}{@{}rrrrrrrrr@{}}
    \toprule
    idx & subject & label & EDA\_mean & EDA\_phasic\_mean & BVP\_std & TEMP\_mean & Resp\_std & net\_acc\_mean \\
    \midrule
    0 & 2 & 1 & 1.398 & 1.829 & 107.679 & 35.817 & 0.011 & 0.026 \\
    1 & 2 & 1 & 1.210 & 2.111 & 118.776 & 35.798 & 0.021 & 0.029 \\
    2 & 2 & 1 & 1.011 & 0.153 & 42.202 & 35.713 & 0.011 & 0.033 \\
    3 & 2 & 1 & 0.775 & 0.178 & 41.619 & 35.701 & 0.003 & 0.018 \\
    \bottomrule
  \end{tabular}
  \caption{Sample windows from \texttt{m14\_merged.csv} (representative subset of columns).}
  \label{tab:samples}
\end{table}

\subsection{Column-by-column explanation}
\renewcommand{\arraystretch}{1.1}
\begin{longtable}{@{}p{0.30\textwidth}p{0.65\textwidth}@{}}
\caption{\label{tab:columns}Column headers in \texttt{m14\_merged.csv} and their meaning.}\\
\toprule
\textbf{Column} & \textbf{Meaning} \\
\midrule
\endfirsthead
\toprule
\textbf{Column} & \textbf{Meaning} \\
\midrule
\endhead
\bottomrule
\endfoot

\texttt{(blank first column)} & CSV index column from export; not a feature. \\

\texttt{EDA\_mean} & Mean of wrist electrodermal activity (EDA) in the window. \\
\texttt{EDA\_std} & Standard deviation of EDA in the window. \\
\texttt{EDA\_min} & Minimum EDA in the window. \\
\texttt{EDA\_max} & Maximum EDA in the window. \\

\texttt{EDA\_phasic\_mean} & Mean of the phasic EDA component in the window. \\
\texttt{EDA\_phasic\_std} & Std of the phasic EDA component in the window. \\
\texttt{EDA\_phasic\_min} & Minimum of the phasic EDA component in the window. \\
\texttt{EDA\_phasic\_max} & Maximum of the phasic EDA component in the window. \\

\texttt{EDA\_smna\_mean} & Mean of the sparse driver component in the window. \\
\texttt{EDA\_smna\_std} & Std of the sparse driver component in the window. \\
\texttt{EDA\_smna\_min} & Minimum of the sparse driver component in the window. \\
\texttt{EDA\_smna\_max} & Maximum of the sparse driver component in the window. \\

\texttt{EDA\_tonic\_mean} & Mean of the tonic EDA component in the window. \\
\texttt{EDA\_tonic\_std} & Std of the tonic EDA component in the window. \\
\texttt{EDA\_tonic\_min} & Minimum of the tonic EDA component in the window. \\
\texttt{EDA\_tonic\_max} & Maximum of the tonic EDA component in the window. \\

\texttt{BVP\_mean} & Mean of wrist blood volume pulse (BVP) in the window. \\
\texttt{BVP\_std} & Std of BVP in the window. \\
\texttt{BVP\_min} & Minimum BVP in the window. \\
\texttt{BVP\_max} & Maximum BVP in the window. \\

\texttt{TEMP\_mean} & Mean wrist skin temperature in the window. \\
\texttt{TEMP\_std} & Std of temperature in the window. \\
\texttt{TEMP\_min} & Minimum temperature in the window. \\
\texttt{TEMP\_max} & Maximum temperature in the window. \\

\texttt{ACC\_x\_mean} & Mean wrist acceleration (x-axis) in the window. \\
\texttt{ACC\_x\_std} & Std wrist acceleration (x-axis) in the window. \\
\texttt{ACC\_x\_min} & Minimum wrist acceleration (x-axis) in the window. \\
\texttt{ACC\_x\_max} & Maximum wrist acceleration (x-axis) in the window. \\

\texttt{ACC\_y\_mean} & Mean wrist acceleration (y-axis) in the window. \\
\texttt{ACC\_y\_std} & Std wrist acceleration (y-axis) in the window. \\
\texttt{ACC\_y\_min} & Minimum wrist acceleration (y-axis) in the window. \\
\texttt{ACC\_y\_max} & Maximum wrist acceleration (y-axis) in the window. \\

\texttt{ACC\_z\_mean} & Mean wrist acceleration (z-axis) in the window. \\
\texttt{ACC\_z\_std} & Std wrist acceleration (z-axis) in the window. \\
\texttt{ACC\_z\_min} & Minimum wrist acceleration (z-axis) in the window. \\
\texttt{ACC\_z\_max} & Maximum wrist acceleration (z-axis) in the window. \\

\texttt{Resp\_mean} & Mean respiration signal in the window. \\
\texttt{Resp\_std} & Std respiration signal in the window. \\
\texttt{Resp\_min} & Minimum respiration signal in the window. \\
\texttt{Resp\_max} & Maximum respiration signal in the window. \\

\texttt{net\_acc\_mean} & Mean acceleration magnitude in the window. \\
\texttt{net\_acc\_std} & Std acceleration magnitude in the window. \\
\texttt{net\_acc\_min} & Minimum acceleration magnitude in the window. \\
\texttt{net\_acc\_max} & Maximum acceleration magnitude in the window. \\

\texttt{BVP\_peak\_freq} & Dominant frequency of BVP in the window. \\
\texttt{TEMP\_slope} & Linear temperature trend across the window. \\

\texttt{subject} & Subject identifier (integer). \\
\texttt{label} & Window label (0=amusement, 1=baseline, 2=stress). \\

\texttt{age} & Subject age. \\
\texttt{height} & Subject height. \\
\texttt{weight} & Subject weight. \\

\texttt{gender\_ female} & One-hot indicator: female. \\
\texttt{gender\_ male} & One-hot indicator: male. \\
\texttt{coffee\_today\_YES} & One-hot indicator: drank coffee today. \\
\texttt{sport\_today\_YES} & One-hot indicator: did sports today. \\
\texttt{smoker\_NO} & One-hot indicator: not a smoker. \\
\texttt{smoker\_YES} & One-hot indicator: smoker. \\
\texttt{feel\_ill\_today\_YES} & One-hot indicator: feels ill today. \\

\end{longtable}

\section{Exploratory data analysis (EDA): what each plot tells us and why it matters}
The EDA is designed to answer three modeling-critical questions:
\begin{enumerate}
  \item \textbf{Is the dataset balanced and comparable across subjects?}
  \item \textbf{Do features separate the classes, and which modalities matter most?}
  \item \textbf{Are there artifacts/outliers or subject-specific effects that will hurt generalization?}
\end{enumerate}

\subsection{Dataset composition and balance}
\begin{figure}[H]
  \centering
  \SubFigIfExists{eda_outputs_m14/01_label_counts.png}{0.49}{Overall class balance.}{fig:label_counts}{}
  \hfill
  \SubFigIfExists{eda_outputs_m14/02_windows_per_subject_stacked.png}{0.49}{Per-subject window counts split by label.}{fig:w_per_subj}{}
  \caption{Composition plots used to assess imbalance and protocol consistency.}
  \label{fig:composition}
\end{figure}

\paragraph{What you can gather.} Class imbalance and subject imbalance guide class weighting, sampling, and which metrics to report (e.g., macro-F1).

\subsection{Univariate feature separation}
\begin{figure}[H]
  \centering
  \SubFigIfExists{eda_outputs_m14/03_box_eda_mean_by_label.png}{0.49}{EDA mean by label.}{fig:eda_mean_box}{}
  \hfill
  \SubFigIfExists{eda_outputs_m14/04_box_eda_phasic_mean_by_label.png}{0.49}{EDA phasic mean by label.}{fig:eda_phasic_box}{}
  \caption{Univariate comparisons to see which features shift across conditions.}
  \label{fig:univariate}
\end{figure}

\paragraph{What you can gather.} Whether stress shifts central tendency, how much overlap exists, and whether outliers/artifacts are present.

\subsection{Movement confounding check}
\FigIfExists{eda_outputs_m14/05_box_net_acc_mean_by_label.png}{0.62}{Acceleration magnitude by label (detects motion confounding).}{fig:motion_conf}

\paragraph{What you can gather.} If motion differs by label, a model may learn artifacts; motivates motion-aware preprocessing and ablations.

\subsection{Correlation and redundancy (original heatmap)}
\FigIfExists{eda_outputs_m14/06_corr_heatmap_subset.png}{0.72}{Correlation heatmap (subset) for redundancy/interaction analysis.}{fig:corr_subset}

\paragraph{What you can gather.} Highly correlated features are redundant; motivates feature selection/regularization.

\subsection{Additional correlation analysis (from \texttt{eda\_m14\_correlation.py})}
\FigIfExists{eda_outputs_m14_correlation/01_corr_overall_spearman.png}{0.78}{Overall correlation (Spearman).}{fig:corr_overall_s}
\FigIfExists{eda_outputs_m14_correlation/02_corr_overall_pearson.png}{0.78}{Overall correlation (Pearson).}{fig:corr_overall_p}

\paragraph{How to interpret.}
\begin{itemize}
  \item \textbf{Spearman} captures monotonic relationships (more robust to heavy tails/outliers).
  \item \textbf{Pearson} captures linear relationships (sensitive to outliers).
  \item Large blocks of high correlation indicate redundant engineered features; consider dropping some or using stronger regularization.
\end{itemize}

\subsubsection*{Condition-specific correlation shifts}
\FigIfExists{eda_outputs_m14_correlation/04_corr_label_1_baseline_spearman.png}{0.78}{Spearman correlation within baseline windows.}{fig:corr_base}
\FigIfExists{eda_outputs_m14_correlation/04_corr_label_2_stress_spearman.png}{0.78}{Spearman correlation within stress windows.}{fig:corr_stress}
\FigIfExists{eda_outputs_m14_correlation/04_corr_label_0_amusement_spearman.png}{0.78}{Spearman correlation within amusement windows.}{fig:corr_amuse}

\paragraph{Why this is useful.} If correlations change by condition (e.g., motion correlates with BVP variability only in stress), it suggests
condition-specific coupling or artifacts, and motivates modality ablations and robustness checks.

\subsection{Low-dimensional projections (PCA)}
\begin{figure}[H]
  \centering
  \SubFigIfExists{eda_outputs_m14/07_pca_scatter_by_label.png}{0.49}{PCA colored by label.}{fig:pca_label}{}
  \hfill
  \SubFigIfExists{eda_outputs_m14/08_pca_scatter_by_subject.png}{0.49}{PCA colored by subject.}{fig:pca_subject}{}
  \caption{PCA projections to compare label structure vs subject structure.}
  \label{fig:pca}
\end{figure}

\paragraph{What you can gather.} If windows cluster by subject more than label, cross-subject generalization is harder and LOSO is required.

\subsection{Window timelines and outlier windows (within-subject)}
\FigIfExists{eda_outputs_m14_timeseries/subject_02_label_timeline.png}{0.92}{Within-subject label segments across window index.}{fig:timeline}
\FigIfExists{eda_outputs_m14_timeseries/subject_02_feature_traces.png}{0.80}{Within-subject feature traces across windows with outliers highlighted.}{fig:traces}

\paragraph{What you can gather.} Drift vs sudden jumps, whether outliers concentrate in certain segments, and whether transitions align with changes.

\subsection{Raw signal inspection with label overlay}
\FigIfExists{eda_outputs_wesad_raw/S02_raw_EDA.png}{0.92}{Raw EDA with protocol labels overlaid.}{fig:raw_eda}
\FigIfExists{eda_outputs_wesad_raw/S02_raw_BVP.png}{0.92}{Raw BVP with protocol labels overlaid (spikes often indicate motion artifacts).}{fig:raw_bvp}
\FigIfExists{eda_outputs_wesad_raw/S02_raw_ACC_mag.png}{0.92}{Raw acceleration magnitude with labels overlaid (validates motion artifacts).}{fig:raw_acc}
\FigIfExists{eda_outputs_wesad_raw/S02_raw_TEMP.png}{0.92}{Raw temperature with labels overlaid (slow trends and drift).}{fig:raw_temp}
\FigIfExists{eda_outputs_wesad_raw/S02_raw_EDA_BVP_combined.png}{0.92}{Combined raw EDA + BVP view with labels overlaid.}{fig:raw_combined}

\paragraph{What you can gather.} Whether BVP spikes align with ACC bursts, segment-level data quality differences, and plausible raw trends.

\subsection{Comparative feature analysis (from \texttt{eda\_m14\_comparative.py})}
\FigIfExists{eda_outputs_m14_comparative/01_bivariate_01_EDA_phasic_mean_vs_BVP_std.png}{0.75}{Example bivariate relationship: EDA\_phasic\_mean vs BVP\_std (colored by label).}{fig:bivar1}
\FigIfExists{eda_outputs_m14_comparative/04_feature_ranking_mutual_info_top20.png}{0.85}{Top features ranked by mutual information with the label.}{fig:mi_rank}
\FigIfExists{eda_outputs_m14_comparative/05_feature_ranking_effectsize_eps2_top20.png}{0.85}{Top features ranked by robust multiclass effect size (epsilon$^2$).}{fig:eps2_rank}

\paragraph{How to interpret.}
\begin{itemize}
  \item Bivariate plots show class separation that may be weak in any single feature but stronger in combinations.
  \item Mutual information highlights nonlinear predictive features; effect size highlights features with strong distribution shifts across classes.
  \item Rankings help choose smaller feature sets and guide sensor ablations (which modalities matter most).
\end{itemize}

\section{Generalization and why LOSO is used}
\subsection{What ``generalization'' means here}
Generalization means: a model trained on some subjects should correctly classify windows from a \textbf{new subject} never seen during training.

\subsection{Why random train/test splits are misleading}
If windows are randomly split, windows from the same subject appear in both train and test.
Because subjects have strong physiological ``signatures'', the model can partially memorize subject-specific patterns, inflating test accuracy.

\subsection{Leave-One-Subject-Out (LOSO)}
LOSO cross-validation holds out one subject for testing and trains on the remaining subjects, repeating this for each subject.
This directly measures subject-independent performance and is the most appropriate estimate of wearable generalization.

\section{How the EDA informs model selection and training}
\begin{itemize}
  \item Use \textbf{LOSO} to estimate cross-subject generalization.
  \item Standardize features and consider per-subject normalization to reduce baseline shifts.
  \item Run modality \textbf{ablations} (motion-only vs physiology-only vs fused) to avoid learning motion artifacts.
  \item Address outliers/artifacts via robust scaling, clipping, or window-quality filtering.
\end{itemize}

\end{document}

\documentclass[11pt]{article}

\usepackage[margin=1in]{geometry}
\usepackage{graphicx}
\usepackage{subcaption}
\usepackage{booktabs}
\usepackage{longtable}
\usepackage{array}
\usepackage[hidelinks]{hyperref}
\usepackage{float}     % enables [H] placement
\usepackage[section]{placeins} % prevents floats crossing sections

\title{WESAD: Data Transformation and Exploratory Data Analysis}
\author{Anishwar Chakraborty}
\date{\today}

\begin{document}
\maketitle

\section{WESAD overview (what it is and what it contains)}
WESAD (Wearable Stress and Affect Detection) is a multimodal wearable-sensing dataset collected in a controlled lab protocol to study how
physiological and movement signals relate to stress and affect.
Each subject experiences multiple \textbf{protocol segments} such as \textbf{baseline} (rest), \textbf{stress} (stress induction),
and \textbf{amusement} (positive affect).

\subsection{What data types are present?}
WESAD is not an image or text dataset. It contains:
\begin{itemize}
  \item \textbf{Time-series sensor signals} from wrist and chest devices (e.g., EDA, BVP/PPG, ACC, temperature, respiration).
  \item \textbf{A synchronized label stream} indicating which protocol segment is active at each moment in time.
  \item \textbf{Subject metadata} (demographics and questionnaire-style yes/no responses) stored in per-subject readme files.
\end{itemize}

\subsection{Why it is useful for stress modeling}
WESAD supports stress detection research because it enables:
\begin{itemize}
  \item \textbf{Multimodal learning}: combining physiological and movement signals.
  \item \textbf{Window-based supervised learning}: converting continuous recordings into labeled examples.
  \item \textbf{Cross-subject generalization}: evaluating whether a model trained on some subjects works on new subjects.
\end{itemize}

\section{Data transformation (implemented in \texttt{data\_wrangling.py})}
Raw WESAD recordings are multi-rate time series (different sensors have different sampling frequencies).
Most EDA and most classifiers are easier to run on a \textbf{fixed-length window} dataset where each row is one labeled example.
The transformation pipeline converts raw signals into a window-level feature table.

\subsection{What the transformation does (conceptually)}
\begin{enumerate}
  \item \textbf{Load raw wrist/chest signals and labels.}
  \item \textbf{Denoise signals} (e.g., smoothing / low-pass filtering) to reduce high-frequency noise and improve robustness.
  \item \textbf{Time-align modalities.} Because sensors are sampled at different rates, signals are aligned onto a consistent timeline.
        Labels are propagated so that every aligned time point (and later every window) has an associated class.
  \item \textbf{EDA decomposition (cvxEDA).} EDA can be decomposed into:
        \textbf{tonic} (slow baseline drift) and \textbf{phasic} (fast skin conductance responses), which can be more informative for stress.
  \item \textbf{Split into protocol segments} (baseline / stress / amusement) using the synchronized label stream.
  \item \textbf{Windowing.} The continuous signal is partitioned into fixed windows (30 seconds), creating a sequence of labeled examples.
  \item \textbf{Feature extraction per window.} For each sensor channel, compute summary statistics (mean/std/min/max).
        Add a few domain-engineered features (e.g., ACC magnitude, BVP dominant frequency, temperature slope).
\end{enumerate}

\subsection{Why summary statistics are used as features}
The window feature columns (e.g., \texttt{EDA\_mean}, \texttt{BVP\_std}, \texttt{TEMP\_max}) are \textbf{summary statistics}
computed over all samples inside one window. The goal is to compress a high-frequency signal into a stable vector that still captures
the window's overall level, variability, and extremes.

\subsubsection*{Mean, median, and mode (what they are)}
\begin{itemize}
  \item \textbf{Mean} is the arithmetic average. It captures the typical level of a signal in the window, but can be sensitive to spikes/outliers.
  \item \textbf{Median} is the middle value after sorting. It is more robust to spikes and is often used when artifacts are common.
        (This project uses mean/std/min/max, but median is a common alternative summary feature.)
  \item \textbf{Mode} is the most frequent value. For continuous physiological signals it is usually less informative unless values are quantized/binned,
        but for discrete variables (e.g., a label sequence) mode corresponds to the most common class.
\end{itemize}

\subsubsection*{Other window statistics used here}
\begin{itemize}
  \item \textbf{Standard deviation (std)} measures variability within the window (often rises with motion artifacts in BVP).
  \item \textbf{Min/max} capture extremes, which can reflect physiological peaks or sensor artifacts.
  \item \textbf{Slope} captures monotonic trends (e.g., temperature drifting).
  \item \textbf{Peak frequency} captures dominant periodic structure (useful for signals like BVP).
  \item \textbf{ACC magnitude} provides a compact movement summary and helps detect confounding between label and motion.
\end{itemize}

\subsection{Sensor modalities and derived features (plain-language)}
Below are short explanations of the main signals and derived features used in this project.

\begin{itemize}
  \item \textbf{EDA (Electrodermal Activity)}: skin conductance (typically in $\mu$S). It reflects sympathetic nervous system activity
        (sweat gland activity) and often increases with stress/arousal, but can drift slowly over time and is sensitive to skin contact.
  \item \textbf{EDA phasic}: the fast-changing component of EDA associated with short skin conductance responses (SCRs). This is commonly
        used as an arousal-related feature because it emphasizes rapid responses rather than slow baseline drift.
  \item \textbf{EDA tonic}: the slow baseline level of EDA (sometimes called SCL). It captures longer-term shifts and drift across the protocol.
  \item \textbf{EDA SMNA / sparse driver}: an auxiliary component from the EDA decomposition that represents estimated ``event-like'' driver activity.
        It can be useful for capturing when rapid EDA responses are being generated.

  \item \textbf{BVP (Blood Volume Pulse / PPG)}: an optical photoplethysmography signal from the wrist. It contains pulsatile information related to
        heart beats (heart rate / variability proxies) but is highly sensitive to motion artifacts and poor sensor contact.
  \item \textbf{BVP peak frequency}: the dominant frequency in the BVP window (from a periodogram). It summarizes periodic structure in the signal
        and can correlate with heart-rate band behavior; large motion artifacts can also distort it.

  \item \textbf{TEMP (Skin temperature)}: wrist skin temperature. It usually changes slowly and can reflect peripheral vasoconstriction/vasodilation
        and environmental effects; it is often more informative through trends than rapid fluctuations.
  \item \textbf{TEMP slope}: the linear trend of temperature within a window, capturing gradual warming/cooling.

  \item \textbf{ACC (Accelerometer)}: 3-axis wrist acceleration. It measures movement/gesture intensity and posture changes.
        ACC features are important both as predictive features and as a way to detect confounding (movement can create artifacts in BVP/EDA).
  \item \textbf{net\_acc (ACC magnitude)}: a single ``overall movement'' value computed from the 3 axes (roughly $\sqrt{x^2+y^2+z^2}$),
        then summarized over the window. This is a compact motion intensity feature.

  \item \textbf{Resp (Respiration)}: breathing-related signal from the chest device. Stress can change breathing rate and variability; respiration is
        often complementary to EDA/BVP for stress detection.
\end{itemize}

\section{Metadata parsing and merging (implemented in \texttt{readme\_parser.py})}
Protocol labels indicate the active condition, but subjects differ in demographic and contextual factors (age, gender, smoking, etc.).
Metadata parsing extracts these fields and encodes them as model-ready columns (e.g., one-hot indicators).

\subsection{Why include metadata?}
Metadata can be used to:
\begin{itemize}
  \item support \textbf{interpretation} (e.g., some physiological baselines differ by subject characteristics),
  \item support \textbf{confound analysis} (checking whether subgroups differ systematically),
  \item optionally provide \textbf{covariates} (with care to avoid leakage in evaluation).
\end{itemize}

\section{Final dataset: \texttt{data/m14\_merged.csv}}
\subsection{What one row represents}
Each row corresponds to one \textbf{30-second window} from one subject, represented by engineered window features plus a class label and metadata.

\subsection{Label meanings}
The labels used in \texttt{m14\_merged.csv} are:
\begin{center}
\begin{tabular}{cl}
\toprule
\textbf{Label value} & \textbf{Meaning} \\
\midrule
0 & amusement \\
1 & baseline \\
2 & stress \\
\bottomrule
\end{tabular}
\end{center}

\subsection{3--4 sample rows (subset of columns)}
The dataset is wide; Table~\ref{tab:samples} shows a small subset of representative columns for 4 windows.

\begin{table}[H]
  \centering
  \small
  \begin{tabular}{@{}rrrrrrrrr@{}}
    \toprule
    idx & subject & label & EDA\_mean & EDA\_phasic\_mean & BVP\_std & TEMP\_mean & Resp\_std & net\_acc\_mean \\
    \midrule
    0 & 2 & 1 & 1.398 & 1.829 & 107.679 & 35.817 & 0.011 & 0.026 \\
    1 & 2 & 1 & 1.210 & 2.111 & 118.776 & 35.798 & 0.021 & 0.029 \\
    2 & 2 & 1 & 1.011 & 0.153 & 42.202 & 35.713 & 0.011 & 0.033 \\
    3 & 2 & 1 & 0.775 & 0.178 & 41.619 & 35.701 & 0.003 & 0.018 \\
    \bottomrule
  \end{tabular}
  \caption{Sample windows from \texttt{m14\_merged.csv} (representative subset of columns).}
  \label{tab:samples}
\end{table}

\subsection{Column-by-column explanation}
Table~\ref{tab:columns} lists every column header and explains what it measures.

\renewcommand{\arraystretch}{1.1}
\begin{longtable}{@{}p{0.30\textwidth}p{0.65\textwidth}@{}}
\caption{\label{tab:columns}Column headers in \texttt{m14\_merged.csv} and their meaning.}\\
\toprule
\textbf{Column} & \textbf{Meaning} \\
\midrule
\endfirsthead
\toprule
\textbf{Column} & \textbf{Meaning} \\
\midrule
\endhead
\bottomrule
\endfoot

\texttt{(blank first column)} & CSV index column from export; not a feature. \\

\texttt{EDA\_mean} & Mean of wrist electrodermal activity (EDA) in the window. \\
\texttt{EDA\_std} & Standard deviation of EDA in the window. \\
\texttt{EDA\_min} & Minimum EDA in the window. \\
\texttt{EDA\_max} & Maximum EDA in the window. \\

\texttt{EDA\_phasic\_mean} & Mean of the phasic EDA component (fast responses) in the window. \\
\texttt{EDA\_phasic\_std} & Std of the phasic EDA component in the window. \\
\texttt{EDA\_phasic\_min} & Minimum of the phasic EDA component in the window. \\
\texttt{EDA\_phasic\_max} & Maximum of the phasic EDA component in the window. \\

\texttt{EDA\_smna\_mean} & Mean of the sparse driver component from EDA decomposition in the window. \\
\texttt{EDA\_smna\_std} & Std of the sparse driver component in the window. \\
\texttt{EDA\_smna\_min} & Minimum of the sparse driver component in the window. \\
\texttt{EDA\_smna\_max} & Maximum of the sparse driver component in the window. \\

\texttt{EDA\_tonic\_mean} & Mean of the tonic EDA component (slow drift / baseline level) in the window. \\
\texttt{EDA\_tonic\_std} & Std of the tonic EDA component in the window. \\
\texttt{EDA\_tonic\_min} & Minimum of the tonic EDA component in the window. \\
\texttt{EDA\_tonic\_max} & Maximum of the tonic EDA component in the window. \\

\texttt{BVP\_mean} & Mean of wrist blood volume pulse (BVP) in the window. \\
\texttt{BVP\_std} & Std of BVP in the window (often increases with motion artifacts). \\
\texttt{BVP\_min} & Minimum BVP in the window. \\
\texttt{BVP\_max} & Maximum BVP in the window. \\

\texttt{TEMP\_mean} & Mean wrist skin temperature in the window. \\
\texttt{TEMP\_std} & Std of temperature in the window. \\
\texttt{TEMP\_min} & Minimum temperature in the window. \\
\texttt{TEMP\_max} & Maximum temperature in the window. \\

\texttt{ACC\_x\_mean} & Mean wrist acceleration (x-axis) in the window. \\
\texttt{ACC\_x\_std} & Std wrist acceleration (x-axis) in the window. \\
\texttt{ACC\_x\_min} & Minimum wrist acceleration (x-axis) in the window. \\
\texttt{ACC\_x\_max} & Maximum wrist acceleration (x-axis) in the window. \\

\texttt{ACC\_y\_mean} & Mean wrist acceleration (y-axis) in the window. \\
\texttt{ACC\_y\_std} & Std wrist acceleration (y-axis) in the window. \\
\texttt{ACC\_y\_min} & Minimum wrist acceleration (y-axis) in the window. \\
\texttt{ACC\_y\_max} & Maximum wrist acceleration (y-axis) in the window. \\

\texttt{ACC\_z\_mean} & Mean wrist acceleration (z-axis) in the window. \\
\texttt{ACC\_z\_std} & Std wrist acceleration (z-axis) in the window. \\
\texttt{ACC\_z\_min} & Minimum wrist acceleration (z-axis) in the window. \\
\texttt{ACC\_z\_max} & Maximum wrist acceleration (z-axis) in the window. \\

\texttt{Resp\_mean} & Mean respiration signal in the window. \\
\texttt{Resp\_std} & Std respiration signal in the window. \\
\texttt{Resp\_min} & Minimum respiration signal in the window. \\
\texttt{Resp\_max} & Maximum respiration signal in the window. \\

\texttt{net\_acc\_mean} & Mean acceleration magnitude (derived from ACC x/y/z) in the window. \\
\texttt{net\_acc\_std} & Std acceleration magnitude in the window. \\
\texttt{net\_acc\_min} & Minimum acceleration magnitude in the window. \\
\texttt{net\_acc\_max} & Maximum acceleration magnitude in the window. \\

\texttt{BVP\_peak\_freq} & Dominant frequency of BVP in the window (periodogram peak). \\
\texttt{TEMP\_slope} & Linear temperature trend across the window. \\

\texttt{subject} & Subject identifier (integer). \\
\texttt{label} & Window label (0=amusement, 1=baseline, 2=stress). \\

\texttt{age} & Subject age. \\
\texttt{height} & Subject height. \\
\texttt{weight} & Subject weight. \\

\texttt{gender\_ female} & One-hot indicator: female. \\
\texttt{gender\_ male} & One-hot indicator: male. \\
\texttt{coffee\_today\_YES} & One-hot indicator: drank coffee today. \\
\texttt{sport\_today\_YES} & One-hot indicator: did sports today. \\
\texttt{smoker\_NO} & One-hot indicator: not a smoker. \\
\texttt{smoker\_YES} & One-hot indicator: smoker. \\
\texttt{feel\_ill\_today\_YES} & One-hot indicator: feels ill today. \\

\end{longtable}

\section{Exploratory data analysis (EDA): what each plot tells us and why it matters}
The EDA is designed to answer three modeling-critical questions:
\begin{enumerate}
  \item \textbf{Is the dataset balanced and comparable across subjects?}
  \item \textbf{Do features separate the classes, and which modalities matter most?}
  \item \textbf{Are there artifacts/outliers or subject-specific effects that will hurt generalization?}
\end{enumerate}

\subsection{Dataset composition and balance}
\begin{figure}[H]
  \centering
  \begin{subfigure}{0.49\linewidth}
    \centering
    \includegraphics[width=\linewidth]{\detokenize{eda_outputs_m14/01_label_counts.png}}
    \caption{Overall class balance.}
  \end{subfigure}
  \begin{subfigure}{0.49\linewidth}
    \centering
    \includegraphics[width=\linewidth]{\detokenize{eda_outputs_m14/02_windows_per_subject_stacked.png}}
    \caption{Per-subject window counts split by label.}
  \end{subfigure}
  \caption{Composition plots used to assess imbalance and protocol consistency.}
  \label{fig:composition}
\end{figure}

\paragraph{What you can gather from these plots.}
\begin{itemize}
  \item Whether the dataset is \textbf{class imbalanced} (one label has far more windows).
  \item Whether each subject contributes a similar number of windows (protocol consistency) or some subjects dominate the dataset.
  \item Whether any class is missing or underrepresented for some subjects (can cause unstable training/evaluation).
\end{itemize}

\paragraph{Why it matters for training.}
These plots guide class-weighting, balanced sampling, and whether to report metrics robust to imbalance (e.g., macro-F1).

\subsection{Univariate feature separation by label}
\begin{figure}[H]
  \centering
  \begin{subfigure}{0.49\linewidth}
    \centering
    \includegraphics[width=\linewidth]{\detokenize{eda_outputs_m14/03_box_eda_mean_by_label.png}}
    \caption{EDA mean by label.}
  \end{subfigure}
  \begin{subfigure}{0.49\linewidth}
    \centering
    \includegraphics[width=\linewidth]{\detokenize{eda_outputs_m14/04_box_eda_phasic_mean_by_label.png}}
    \caption{EDA phasic mean by label.}
  \end{subfigure}
  \caption{Univariate comparisons to see which physiological features shift across conditions.}
  \label{fig:univariate}
\end{figure}

\paragraph{What you can gather from these plots.}
\begin{itemize}
  \item Whether stress tends to have \textbf{higher central tendency} than baseline/amusement.
  \item How much \textbf{overlap} exists between classes (heavy overlap suggests single features are not sufficient).
  \item Presence of \textbf{outliers} that may correspond to artifacts or rare physiological events.
\end{itemize}

\paragraph{Why it matters for training.}
If a single modality shows strong separation, simpler models may perform well; if separation is weak, multimodal fusion helps.

\subsection{Movement confounding check}
\begin{figure}[H]
  \centering
  \includegraphics[width=0.62\linewidth]{\detokenize{eda_outputs_m14/05_box_net_acc_mean_by_label.png}}
  \caption{Acceleration magnitude by label (used to detect motion confounding).}
  \label{fig:motion_conf}
\end{figure}

\paragraph{What you can gather from this plot.}
\begin{itemize}
  \item Whether one label coincides with systematically higher movement.
  \item Whether there is a potential \textbf{confound}: ``stress'' $\approx$ ``movement'' rather than ``stress'' $\approx$ ``physiology''.
\end{itemize}

\paragraph{Why it matters for training.}
If movement differs by label, a classifier may learn motion artifacts (especially in wrist BVP).

\subsection{Correlation and redundancy}
\begin{figure}[H]
  \centering
  \includegraphics[width=0.72\linewidth]{\detokenize{eda_outputs_m14/06_corr_heatmap_subset.png}}
  \caption{Correlation heatmap (subset) for redundancy/interaction analysis.}
  \label{fig:corr}
\end{figure}

\paragraph{What you can gather from this plot.}
\begin{itemize}
  \item Which features are highly redundant (very high correlation).
  \item Which modalities move together (e.g., motion correlating with BVP variability).
\end{itemize}

\paragraph{Why it matters for training.}
Redundancy motivates feature selection, dimensionality reduction, and/or stronger regularization.

\subsection{Low-dimensional projections (PCA)}
\begin{figure}[H]
  \centering
  \begin{subfigure}{0.49\linewidth}
    \centering
    \includegraphics[width=\linewidth]{\detokenize{eda_outputs_m14/07_pca_scatter_by_label.png}}
    \caption{PCA colored by label.}
  \end{subfigure}
  \begin{subfigure}{0.49\linewidth}
    \centering
    \includegraphics[width=\linewidth]{\detokenize{eda_outputs_m14/08_pca_scatter_by_subject.png}}
    \caption{PCA colored by subject.}
  \end{subfigure}
  \caption{PCA projections to compare label structure vs subject structure.}
  \label{fig:pca}
\end{figure}

\paragraph{What you can gather from these plots.}
\begin{itemize}
  \item Whether windows cluster primarily by \textbf{label} or by \textbf{subject}.
  \item Whether there are extreme outliers that might distort training.
\end{itemize}

\paragraph{Why it matters for training.}
Subject clustering is a sign of cross-subject distribution shift and motivates LOSO evaluation and normalization strategies.

\subsection{Window timelines and outlier windows (within-subject)}
\begin{figure}[H]
  \centering
  \includegraphics[width=0.92\linewidth]{\detokenize{eda_outputs_m14_timeseries/subject_02_label_timeline.png}}
  \caption{Within-subject label segments across window index.}
  \label{fig:timeline}
\end{figure}

\begin{figure}[H]
  \centering
  \includegraphics[width=0.78\linewidth]{\detokenize{eda_outputs_m14_timeseries/subject_02_feature_traces.png}}
  \caption{Within-subject feature traces across windows with outliers highlighted.}
  \label{fig:traces}
\end{figure}

\paragraph{What you can gather from these plots.}
\begin{itemize}
  \item Whether features drift gradually or jump suddenly (possible artifacts).
  \item Whether outliers concentrate in particular protocol segments.
  \item Whether label transitions align with abrupt physiological changes (sanity check).
\end{itemize}

\paragraph{Why it matters for training.}
Outlier windows can dominate training and degrade generalization; these plots motivate robust scaling and window filtering.

\subsection{Raw signal inspection with label overlay}
\begin{figure}[H]
  \centering
  \includegraphics[width=0.92\linewidth]{\detokenize{eda_outputs_wesad_raw/S02_raw_EDA.png}}
  \caption{Raw EDA with protocol labels overlaid.}
  \label{fig:raw_eda}
\end{figure}

\begin{figure}[H]
  \centering
  \includegraphics[width=0.92\linewidth]{\detokenize{eda_outputs_wesad_raw/S02_raw_BVP.png}}
  \caption{Raw BVP with protocol labels overlaid (spikes often indicate motion artifacts).}
  \label{fig:raw_bvp}
\end{figure}

\begin{figure}[H]
  \centering
  \includegraphics[width=0.92\linewidth]{\detokenize{eda_outputs_wesad_raw/S02_raw_ACC_mag.png}}
  \caption{Raw acceleration magnitude with labels overlaid (used to validate motion artifacts).}
  \label{fig:raw_acc}
\end{figure}

\begin{figure}[H]
  \centering
  \includegraphics[width=0.92\linewidth]{\detokenize{eda_outputs_wesad_raw/S02_raw_TEMP.png}}
  \caption{Raw temperature with labels overlaid (slow trends and drift).}
  \label{fig:raw_temp}
\end{figure}

\begin{figure}[H]
  \centering
  \includegraphics[width=0.92\linewidth]{\detokenize{eda_outputs_wesad_raw/S02_raw_EDA_BVP_combined.png}}
  \caption{Combined raw EDA + BVP view with labels overlaid.}
  \label{fig:raw_combined}
\end{figure}

\paragraph{What you can gather from these plots.}
\begin{itemize}
  \item Whether BVP spikes co-occur with ACC bursts (motion artifacts).
  \item Whether one segment is systematically noisier than others (data quality differs by condition).
  \item Whether raw trends (EDA, TEMP) plausibly change across the protocol segments.
\end{itemize}

\paragraph{Why it matters for training.}
Raw inspection prevents misinterpretation of window features and supports artifact-aware preprocessing and window-quality rules.

\section{Generalization and why LOSO is used}
\subsection{What ``generalization'' means here}
In this project, \textbf{generalization} means: a model trained on some subjects should correctly classify windows from a \textbf{new subject}
it has never seen during training.

\subsection{Why random train/test splits are misleading}
If windows are randomly split across train and test, windows from the \emph{same subject} appear in both sets.
Because each subject has a strong physiological ``signature'', the model can partially memorize subject-specific patterns.
This yields overly optimistic test performance that does not reflect real-world performance on new users.

\subsection{Leave-One-Subject-Out (LOSO)}
LOSO cross-validation evaluates generalization by holding out one subject as the test set and training on all remaining subjects.
This is repeated for each subject; the final performance is the average (and often the variability) across held-out subjects.

\paragraph{Why LOSO is appropriate for WESAD.}
\begin{itemize}
  \item It directly measures subject-independent performance (the realistic wearable deployment setting).
  \item It tests the model under distribution shift (different baselines across people).
  \item It discourages leakage and inflated accuracy caused by within-subject correlations.
\end{itemize}

\section{How the EDA informs model selection and training}
The transformation and EDA support the following decisions:
\begin{itemize}
  \item \textbf{Validation strategy}: use LOSO to estimate cross-subject generalization.
  \item \textbf{Scaling and normalization}: standardize features (and consider per-subject normalization) to reduce baseline shifts.
  \item \textbf{Confound-aware modeling}: compare motion-only vs physiology-only vs fused models to ensure the model learns stress physiology.
  \item \textbf{Robustness}: address outliers/artifacts via robust scaling, clipping, and/or removal of low-quality windows.
  \item \textbf{Feature/model choice}: if univariate plots show limited separation, prioritize multimodal fusion and stronger regularization.
\end{itemize}

\end{document}

\documentclass[11pt]{article}

\usepackage[margin=1in]{geometry}
\usepackage{graphicx}
\usepackage{subcaption}
\usepackage{booktabs}
\usepackage{longtable}
\usepackage{array}
\usepackage[hidelinks]{hyperref}

\title{WESAD: Data Transformation and Exploratory Data Analysis}
\author{Anishwar Chakraborty}
\date{\today}

\begin{document}
\maketitle

\section{WESAD overview (what it is and what it contains)}
WESAD (Wearable Stress and Affect Detection) is a multimodal wearable-sensing dataset collected in a controlled lab protocol to study how
physiological and movement signals relate to stress and affect.
Each subject experiences multiple \textbf{protocol segments} such as \textbf{baseline} (rest), \textbf{stress} (stress induction),
and \textbf{amusement} (positive affect).

\subsection{What data types are present?}
WESAD is not an image or text dataset. It contains:
\begin{itemize}
  \item \textbf{Time-series sensor signals} from wrist and chest devices (e.g., EDA, BVP/PPG, ACC, temperature, respiration).
  \item \textbf{A synchronized label stream} indicating which protocol segment is active at each moment in time.
  \item \textbf{Subject metadata} (demographics and questionnaire-style yes/no responses) stored in per-subject readme files.
\end{itemize}

\subsection{Why it is useful for stress modeling}
WESAD supports stress detection research because it enables:
\begin{itemize}
  \item \textbf{Multimodal learning}: combining physiological and movement signals.
  \item \textbf{Window-based supervised learning}: converting continuous recordings into labeled examples.
  \item \textbf{Cross-subject generalization}: evaluating whether a model trained on some subjects works on new subjects.
\end{itemize}

\section{Data transformation (implemented in \texttt{data\_wrangling.py})}
Raw WESAD recordings are multi-rate time series (different sensors have different sampling frequencies).
Most EDA and most classifiers are easier to run on a \textbf{fixed-length window} dataset where each row is one labeled example.
The transformation pipeline converts raw signals into a window-level feature table.

\subsection{What the transformation does (conceptually)}
\begin{enumerate}
  \item \textbf{Load raw wrist/chest signals and labels.}
  \item \textbf{Denoise signals} (e.g., smoothing / low-pass filtering) to reduce high-frequency noise and improve robustness.
  \item \textbf{Time-align modalities.} Because sensors are sampled at different rates, signals are aligned onto a consistent timeline.
        Labels are propagated so that every aligned time point (and later every window) has an associated class.
  \item \textbf{EDA decomposition (cvxEDA).} EDA can be decomposed into:
        \textbf{tonic} (slow baseline drift) and \textbf{phasic} (fast skin conductance responses), which can be more informative for stress.
  \item \textbf{Split into protocol segments} (baseline / stress / amusement) using the synchronized label stream.
  \item \textbf{Windowing.} The continuous signal is partitioned into fixed windows (30 seconds), creating a sequence of labeled examples.
  \item \textbf{Feature extraction per window.} For each sensor channel, compute summary statistics (mean/std/min/max).
        Add a few domain-engineered features (e.g., ACC magnitude, BVP dominant frequency, temperature slope).
\end{enumerate}

\subsection{Why summary statistics are used as features}
The window feature columns (e.g., \texttt{EDA\_mean}, \texttt{BVP\_std}, \texttt{TEMP\_max}) are \textbf{summary statistics}
computed over all samples inside one window. The goal is to compress a high-frequency signal into a stable vector that still captures
the window's overall level, variability, and extremes.

\subsubsection*{Mean, median, and mode (what they are)}
\begin{itemize}
  \item \textbf{Mean} is the arithmetic average. It captures the typical level of a signal in the window, but can be sensitive to spikes/outliers.
  \item \textbf{Median} is the middle value after sorting. It is more robust to spikes and is often used when artifacts are common.
        (This project uses mean/std/min/max, but median is a common alternative summary feature.)
  \item \textbf{Mode} is the most frequent value. For continuous physiological signals it is usually less informative unless values are quantized/binned,
        but for discrete variables (e.g., a label sequence) mode corresponds to the most common class.
\end{itemize}

\subsubsection*{Other window statistics used here}
\begin{itemize}
  \item \textbf{Standard deviation (std)} measures variability within the window (often rises with motion artifacts in BVP).
  \item \textbf{Min/max} capture extremes, which can reflect physiological peaks or sensor artifacts.
  \item \textbf{Slope} captures monotonic trends (e.g., temperature drifting).
  \item \textbf{Peak frequency} captures dominant periodic structure (useful for signals like BVP).
  \item \textbf{ACC magnitude} provides a compact movement summary and helps detect confounding between label and motion.
\end{itemize}

\section{Metadata parsing and merging (implemented in \texttt{readme\_parser.py})}
Protocol labels indicate the active condition, but subjects differ in demographic and contextual factors (age, gender, smoking, etc.).
Metadata parsing extracts these fields and encodes them as model-ready columns (e.g., one-hot indicators).

\subsection{Why include metadata?}
Metadata can be used to:
\begin{itemize}
  \item support \textbf{interpretation} (e.g., some physiological baselines differ by subject characteristics),
  \item support \textbf{confound analysis} (checking whether subgroups differ systematically),
  \item optionally provide \textbf{covariates} (with care to avoid leakage in evaluation).
\end{itemize}

\section{Final dataset: \texttt{data/m14\_merged.csv}}
\subsection{What one row represents}
Each row corresponds to one \textbf{30-second window} from one subject, represented by engineered window features plus a class label and metadata.

\subsection{Label meanings}
The labels used in \texttt{m14\_merged.csv} are:
\begin{center}
\begin{tabular}{cl}
\toprule
\textbf{Label value} & \textbf{Meaning} \\
\midrule
0 & amusement \\
1 & baseline \\
2 & stress \\
\bottomrule
\end{tabular}
\end{center}

\subsection{3--4 sample rows (subset of columns)}
The dataset is wide; Table~\ref{tab:samples} shows a small subset of representative columns for 4 windows.

\begin{table}[!ht]
  \centering
  \small
  \begin{tabular}{@{}rrrrrrrrr@{}}
    \toprule
    idx & subject & label & EDA\_mean & EDA\_phasic\_mean & BVP\_std & TEMP\_mean & Resp\_std & net\_acc\_mean \\
    \midrule
    0 & 2 & 1 & 1.398 & 1.829 & 107.679 & 35.817 & 0.011 & 0.026 \\
    1 & 2 & 1 & 1.210 & 2.111 & 118.776 & 35.798 & 0.021 & 0.029 \\
    2 & 2 & 1 & 1.011 & 0.153 & 42.202 & 35.713 & 0.011 & 0.033 \\
    3 & 2 & 1 & 0.775 & 0.178 & 41.619 & 35.701 & 0.003 & 0.018 \\
    \bottomrule
  \end{tabular}
  \caption{Sample windows from \texttt{m14\_merged.csv} (representative subset of columns).}
  \label{tab:samples}
\end{table}

\subsection{Column-by-column explanation}
Table~\ref{tab:columns} lists every column header and explains what it measures.

\renewcommand{\arraystretch}{1.1}
\begin{longtable}{@{}p{0.30\textwidth}p{0.65\textwidth}@{}}
\caption{\label{tab:columns}Column headers in \texttt{m14\_merged.csv} and their meaning.}\\
\toprule
\textbf{Column} & \textbf{Meaning} \\
\midrule
\endfirsthead
\toprule
\textbf{Column} & \textbf{Meaning} \\
\midrule
\endhead
\bottomrule
\endfoot

\texttt{(blank first column)} & CSV index column from export; not a feature. \\

\texttt{EDA\_mean} & Mean of wrist electrodermal activity (EDA) in the window. \\
\texttt{EDA\_std} & Standard deviation of EDA in the window. \\
\texttt{EDA\_min} & Minimum EDA in the window. \\
\texttt{EDA\_max} & Maximum EDA in the window. \\

\texttt{EDA\_phasic\_mean} & Mean of the phasic EDA component (fast responses) in the window. \\
\texttt{EDA\_phasic\_std} & Std of the phasic EDA component in the window. \\
\texttt{EDA\_phasic\_min} & Minimum of the phasic EDA component in the window. \\
\texttt{EDA\_phasic\_max} & Maximum of the phasic EDA component in the window. \\

\texttt{EDA\_smna\_mean} & Mean of the sparse driver component from EDA decomposition in the window. \\
\texttt{EDA\_smna\_std} & Std of the sparse driver component in the window. \\
\texttt{EDA\_smna\_min} & Minimum of the sparse driver component in the window. \\
\texttt{EDA\_smna\_max} & Maximum of the sparse driver component in the window. \\

\texttt{EDA\_tonic\_mean} & Mean of the tonic EDA component (slow drift / baseline level) in the window. \\
\texttt{EDA\_tonic\_std} & Std of the tonic EDA component in the window. \\
\texttt{EDA\_tonic\_min} & Minimum of the tonic EDA component in the window. \\
\texttt{EDA\_tonic\_max} & Maximum of the tonic EDA component in the window. \\

\texttt{BVP\_mean} & Mean of wrist blood volume pulse (BVP) in the window. \\
\texttt{BVP\_std} & Std of BVP in the window (often increases with motion artifacts). \\
\texttt{BVP\_min} & Minimum BVP in the window. \\
\texttt{BVP\_max} & Maximum BVP in the window. \\

\texttt{TEMP\_mean} & Mean wrist skin temperature in the window. \\
\texttt{TEMP\_std} & Std of temperature in the window. \\
\texttt{TEMP\_min} & Minimum temperature in the window. \\
\texttt{TEMP\_max} & Maximum temperature in the window. \\

\texttt{ACC\_x\_mean} & Mean wrist acceleration (x-axis) in the window. \\
\texttt{ACC\_x\_std} & Std wrist acceleration (x-axis) in the window. \\
\texttt{ACC\_x\_min} & Minimum wrist acceleration (x-axis) in the window. \\
\texttt{ACC\_x\_max} & Maximum wrist acceleration (x-axis) in the window. \\

\texttt{ACC\_y\_mean} & Mean wrist acceleration (y-axis) in the window. \\
\texttt{ACC\_y\_std} & Std wrist acceleration (y-axis) in the window. \\
\texttt{ACC\_y\_min} & Minimum wrist acceleration (y-axis) in the window. \\
\texttt{ACC\_y\_max} & Maximum wrist acceleration (y-axis) in the window. \\

\texttt{ACC\_z\_mean} & Mean wrist acceleration (z-axis) in the window. \\
\texttt{ACC\_z\_std} & Std wrist acceleration (z-axis) in the window. \\
\texttt{ACC\_z\_min} & Minimum wrist acceleration (z-axis) in the window. \\
\texttt{ACC\_z\_max} & Maximum wrist acceleration (z-axis) in the window. \\

\texttt{Resp\_mean} & Mean respiration signal in the window. \\
\texttt{Resp\_std} & Std respiration signal in the window. \\
\texttt{Resp\_min} & Minimum respiration signal in the window. \\
\texttt{Resp\_max} & Maximum respiration signal in the window. \\

\texttt{net\_acc\_mean} & Mean acceleration magnitude (derived from ACC x/y/z) in the window. \\
\texttt{net\_acc\_std} & Std acceleration magnitude in the window. \\
\texttt{net\_acc\_min} & Minimum acceleration magnitude in the window. \\
\texttt{net\_acc\_max} & Maximum acceleration magnitude in the window. \\

\texttt{BVP\_peak\_freq} & Dominant frequency of BVP in the window (periodogram peak). \\
\texttt{TEMP\_slope} & Linear temperature trend across the window. \\

\texttt{subject} & Subject identifier (integer). \\
\texttt{label} & Window label (0=amusement, 1=baseline, 2=stress). \\

\texttt{age} & Subject age. \\
\texttt{height} & Subject height. \\
\texttt{weight} & Subject weight. \\

\texttt{gender\_ female} & One-hot indicator: female. \\
\texttt{gender\_ male} & One-hot indicator: male. \\
\texttt{coffee\_today\_YES} & One-hot indicator: drank coffee today. \\
\texttt{sport\_today\_YES} & One-hot indicator: did sports today. \\
\texttt{smoker\_NO} & One-hot indicator: not a smoker. \\
\texttt{smoker\_YES} & One-hot indicator: smoker. \\
\texttt{feel\_ill\_today\_YES} & One-hot indicator: feels ill today. \\

\end{longtable}

\section{Exploratory data analysis (EDA): what each plot tells us and why it matters}
The EDA is designed to answer three modeling-critical questions:
\begin{enumerate}
  \item \textbf{Is the dataset balanced and comparable across subjects?}
  \item \textbf{Do features separate the classes, and which modalities matter most?}
  \item \textbf{Are there artifacts/outliers or subject-specific effects that will hurt generalization?}
\end{enumerate}

\subsection{Dataset composition and balance}
\begin{figure}[!ht]
  \centering
  \begin{subfigure}{0.49\linewidth}
    \centering
    \includegraphics[width=\linewidth]{\detokenize{eda_outputs_m14/01_label_counts.png}}
    \caption{Overall class balance.}
  \end{subfigure}
  \begin{subfigure}{0.49\linewidth}
    \centering
    \includegraphics[width=\linewidth]{\detokenize{eda_outputs_m14/02_windows_per_subject_stacked.png}}
    \caption{Per-subject window counts split by label.}
  \end{subfigure}
  \caption{Composition plots used to assess imbalance and protocol consistency.}
  \label{fig:composition}
\end{figure}

\paragraph{What you can gather from these plots.}
\begin{itemize}
  \item Whether the dataset is \textbf{class imbalanced} (one label has far more windows).
  \item Whether each subject contributes a similar number of windows (protocol consistency) or some subjects dominate the dataset.
  \item Whether any class is missing or underrepresented for some subjects (can cause unstable training/evaluation).
\end{itemize}

\paragraph{Why it matters for training.}
These plots guide class-weighting, balanced sampling, and whether to report metrics robust to imbalance (e.g., macro-F1).

\subsection{Univariate feature separation by label}
\begin{figure}[!ht]
  \centering
  \begin{subfigure}{0.49\linewidth}
    \centering
    \includegraphics[width=\linewidth]{\detokenize{eda_outputs_m14/03_box_eda_mean_by_label.png}}
    \caption{EDA mean by label.}
  \end{subfigure}
  \begin{subfigure}{0.49\linewidth}
    \centering
    \includegraphics[width=\linewidth]{\detokenize{eda_outputs_m14/04_box_eda_phasic_mean_by_label.png}}
    \caption{EDA phasic mean by label.}
  \end{subfigure}
  \caption{Univariate comparisons to see which physiological features shift across conditions.}
  \label{fig:univariate}
\end{figure}

\paragraph{What you can gather from these plots.}
\begin{itemize}
  \item Whether stress tends to have \textbf{higher central tendency} (e.g., EDA\_mean) than baseline/amusement.
  \item How much \textbf{overlap} exists between classes (if boxes/whiskers overlap heavily, single features are not sufficient).
  \item Presence of \textbf{outliers} that may correspond to artifacts or rare physiological events.
\end{itemize}

\paragraph{Why it matters for training.}
If a single modality shows strong separation, simpler models may perform well.
If separation is weak, multimodal fusion and stronger regularization become more important.

\subsection{Movement confounding check}
\begin{figure}[!ht]
  \centering
  \includegraphics[width=0.62\linewidth]{\detokenize{eda_outputs_m14/05_box_net_acc_mean_by_label.png}}
  \caption{Acceleration magnitude by label (used to detect motion confounding).}
  \label{fig:motion_conf}
\end{figure}

\paragraph{What you can gather from this plot.}
\begin{itemize}
  \item Whether one label coincides with systematically higher movement.
  \item Whether the dataset risks a \textbf{confound}: ``stress'' $\approx$ ``movement'' rather than ``stress'' $\approx$ ``physiology''.
\end{itemize}

\paragraph{Why it matters for training.}
If movement differs by label, a classifier may learn motion artifacts (especially in wrist BVP).
This motivates ablations (motion-only vs physiology-only vs fused) and window-quality filtering.

\subsection{Correlation and redundancy}
\begin{figure}[!ht]
  \centering
  \includegraphics[width=0.72\linewidth]{\detokenize{eda_outputs_m14/06_corr_heatmap_subset.png}}
  \caption{Correlation heatmap (subset) for redundancy/interaction analysis.}
  \label{fig:corr}
\end{figure}

\paragraph{What you can gather from this plot.}
\begin{itemize}
  \item Which features are highly redundant (very high correlation).
  \item Which modalities move together (e.g., motion correlating with BVP variability).
\end{itemize}

\paragraph{Why it matters for training.}
Redundancy can inflate model complexity and overfitting; it motivates feature selection, dimensionality reduction, and regularization.

\subsection{Low-dimensional projections (PCA)}
\begin{figure}[!ht]
  \centering
  \begin{subfigure}{0.49\linewidth}
    \centering
    \includegraphics[width=\linewidth]{\detokenize{eda_outputs_m14/07_pca_scatter_by_label.png}}
    \caption{PCA colored by label.}
  \end{subfigure}
  \begin{subfigure}{0.49\linewidth}
    \centering
    \includegraphics[width=\linewidth]{\detokenize{eda_outputs_m14/08_pca_scatter_by_subject.png}}
    \caption{PCA colored by subject.}
  \end{subfigure}
  \caption{PCA projections to compare label structure vs subject structure.}
  \label{fig:pca}
\end{figure}

\paragraph{What you can gather from these plots.}
\begin{itemize}
  \item If the points cluster primarily by \textbf{label}, the features carry condition information.
  \item If the points cluster primarily by \textbf{subject}, subject-specific baselines dominate (distribution shift).
  \item Extreme points can indicate outlier windows that may distort training.
\end{itemize}

\paragraph{Why it matters for training.}
Strong subject clustering motivates subject-independent evaluation (LOSO) and normalization strategies.

\subsection{Window timelines and outlier windows (within-subject)}
\begin{figure}[!ht]
  \centering
  \includegraphics[width=0.92\linewidth]{\detokenize{eda_outputs_m14_timeseries/subject_02_label_timeline.png}}
  \caption{Within-subject label segments across window index.}
  \label{fig:timeline}
\end{figure}

\begin{figure}[!ht]
  \centering
  \includegraphics[width=0.78\linewidth]{\detokenize{eda_outputs_m14_timeseries/subject_02_feature_traces.png}}
  \caption{Within-subject feature traces across windows with outliers highlighted.}
  \label{fig:traces}
\end{figure}

\paragraph{What you can gather from these plots.}
\begin{itemize}
  \item Whether features drift gradually (baseline changes) or jump suddenly (possible artifacts).
  \item Whether outliers concentrate in particular protocol segments.
  \item Whether label transitions align with abrupt physiological changes (sanity check).
\end{itemize}

\paragraph{Why it matters for training.}
Outlier windows can dominate gradients and degrade generalization; these plots motivate robust scaling and window filtering.

\subsection{Raw signal inspection with label overlay}
\begin{figure}[!ht]
  \centering
  \includegraphics[width=0.92\linewidth]{\detokenize{eda_outputs_wesad_raw/S02_raw_EDA.png}}
  \caption{Raw EDA with protocol labels overlaid.}
  \label{fig:raw_eda}
\end{figure}

\begin{figure}[!ht]
  \centering
  \includegraphics[width=0.92\linewidth]{\detokenize{eda_outputs_wesad_raw/S02_raw_BVP.png}}
  \caption{Raw BVP with protocol labels overlaid (spikes often indicate motion artifacts).}
  \label{fig:raw_bvp}
\end{figure}

\begin{figure}[!ht]
  \centering
  \includegraphics[width=0.92\linewidth]{\detokenize{eda_outputs_wesad_raw/S02_raw_ACC_mag.png}}
  \caption{Raw acceleration magnitude with labels overlaid (used to validate motion artifacts).}
  \label{fig:raw_acc}
\end{figure}

\begin{figure}[!ht]
  \centering
  \includegraphics[width=0.92\linewidth]{\detokenize{eda_outputs_wesad_raw/S02_raw_TEMP.png}}
  \caption{Raw temperature with labels overlaid (slow trends and drift).}
  \label{fig:raw_temp}
\end{figure}

\begin{figure}[!ht]
  \centering
  \includegraphics[width=0.92\linewidth]{\detokenize{eda_outputs_wesad_raw/S02_raw_EDA_BVP_combined.png}}
  \caption{Combined raw EDA + BVP view with labels overlaid.}
  \label{fig:raw_combined}
\end{figure}

\paragraph{What you can gather from these plots.}
\begin{itemize}
  \item Whether BVP ``spikes'' co-occur with ACC bursts (motion artifacts).
  \item Whether one segment is systematically noisier than others (data quality differs by condition).
  \item Whether raw trends (EDA, TEMP) plausibly change across the protocol segments.
\end{itemize}

\paragraph{Why it matters for training.}
Raw inspection prevents misinterpretation of window features and supports artifact-aware preprocessing and window-quality rules.

\section{Generalization and why LOSO is used}
\subsection{What ``generalization'' means here}
In this project, \textbf{generalization} means: a model trained on some subjects should correctly classify windows from a \textbf{new subject}
it has never seen during training.

\subsection{Why random train/test splits are misleading}
If windows are randomly split across train and test, windows from the \emph{same subject} appear in both sets.
Because a subject's physiology and sensor placement create a strong ``signature'', the model can partially memorize subject-specific patterns.
This yields overly optimistic test performance that does not reflect real-world performance on new users.

\subsection{Leave-One-Subject-Out (LOSO)}
LOSO cross-validation evaluates generalization by holding out one subject as the test set and training on all remaining subjects.
This is repeated for each subject, producing one score per held-out subject; the final reported performance is the average (and often the variance)
across subjects.

\paragraph{Why LOSO is appropriate for WESAD.}
\begin{itemize}
  \item It directly measures subject-independent performance (the most realistic wearable deployment setting).
  \item It stresses the model under distribution shift (different baselines across people).
  \item It discourages leakage and inflated accuracy caused by within-subject correlations.
\end{itemize}

\section{How the EDA informs model selection and training}
The transformation and EDA support the following decisions:
\begin{itemize}
  \item \textbf{Validation strategy}: use LOSO to estimate cross-subject generalization.
  \item \textbf{Scaling and normalization}: standardize features (and consider per-subject normalization) to reduce baseline shifts.
  \item \textbf{Confound-aware modeling}: compare motion-only vs physiology-only vs fused models to ensure the model learns stress physiology.
  \item \textbf{Robustness}: address outliers/artifacts via robust scaling, clipping, and/or removal of low-quality windows.
  \item \textbf{Feature/model choice}: if univariate plots show limited separation, prioritize multimodal fusion and stronger regularization.
\end{itemize}

\end{document}

\documentclass[11pt]{article}

\usepackage[margin=1in]{geometry}
\usepackage{graphicx}
\usepackage{subcaption}
\usepackage{booktabs}
\usepackage{longtable}
\usepackage{array}
\usepackage[hidelinks]{hyperref}

\title{WESAD: Data Transformation and Exploratory Data Analysis}
\author{Anishwar Chakraborty}
\date{\today}

\begin{document}
\maketitle

\section{WESAD overview (what it is and what it contains)}
WESAD (Wearable Stress and Affect Detection) is a multimodal wearable-sensing dataset collected in a controlled laboratory protocol
to study how physiological and motion signals relate to stress and affect.
Each subject completes multiple \textbf{protocol segments} such as \textbf{baseline} (rest), \textbf{stress} (stress induction),
and \textbf{amusement} (positive affect).

\subsection{What data types are present?}
WESAD is not an image or text dataset. It primarily contains:
\begin{itemize}
  \item \textbf{Time-series sensor signals} from wrist and chest devices (e.g., EDA, BVP/PPG, ACC, temperature, respiration).
  \item \textbf{A synchronized label stream} indicating which protocol segment is active at each moment in time.
  \item \textbf{Subject metadata} (demographics and questionnaire-style yes/no responses) stored in per-subject readme files.
\end{itemize}

\subsection{Why it is useful for stress modeling}
WESAD supports stress detection research because it enables:
\begin{itemize}
  \item \textbf{Multimodal learning}: combining physiological and movement signals.
  \item \textbf{Window-based supervised learning}: converting continuous recordings into labeled examples.
  \item \textbf{Cross-subject generalization}: evaluating whether a model trained on some subjects generalizes to new subjects.
\end{itemize}

\section{Data transformation pipeline (implemented in \texttt{data\_wrangling.py})}
Raw WESAD recordings are high-frequency, multi-rate time series. Most models (and most EDA) are easier when the data is transformed into a
\textbf{table of fixed-length windows} where each row is one labeled example. The transformation pipeline accomplishes that.

\subsection{What the transformation does (conceptually)}
\begin{enumerate}
  \item \textbf{Load raw wrist/chest signals and labels.}
  \item \textbf{Basic denoising} (e.g., smoothing or low-pass filtering) to reduce sensor noise that does not carry useful information.
  \item \textbf{Time alignment across modalities.}
  Signals are sampled at different rates (for example, EDA is much slower than BVP or ACC), so the pipeline aligns them onto a common timeline.
  The label stream is then propagated so every time point (and later every window) receives a label.
  \item \textbf{EDA decomposition (cvxEDA).}
  Electrodermal activity contains both slow drift and rapid responses. Decomposing EDA into \textbf{tonic} and \textbf{phasic} components helps
  separate background level from fast arousal-related changes.
  \item \textbf{Segmenting and windowing.}
  The aligned recording is split into contiguous windows of fixed duration (30 seconds). Each window becomes one training example.
  \item \textbf{Feature extraction per window.}
  For each sensor channel in a window, the pipeline computes summary statistics such as mean, standard deviation, minimum, and maximum.
  Additional engineered features capture domain signals such as:
  \begin{itemize}
    \item \textbf{Motion magnitude} from ACC (useful both as a predictor and as a confound indicator).
    \item \textbf{Dominant frequency} in BVP (a proxy for periodic structure; related to heart-rate band behavior).
    \item \textbf{Temperature slope} (slow trends over the window).
  \end{itemize}
\end{enumerate}

\subsection{Why these steps are used (the modeling motivation)}
\begin{itemize}
  \item \textbf{Fixed windows} turn continuous recordings into a standard supervised learning format.
  \item \textbf{Summary statistics} provide strong baselines and make models less sensitive to raw-signal sampling differences.
  \item \textbf{EDA tonic/phasic features} are often more discriminative than raw EDA alone because stress can affect rapid skin conductance responses.
  \item \textbf{Motion features} are essential because wrist BVP and EDA can be corrupted by motion; without motion-aware analysis a model may learn artifacts.
\end{itemize}

\section{Metadata parsing and merging (implemented in \texttt{readme\_parser.py})}
The protocol labels describe \emph{what segment is happening}, but subjects also differ in demographic and contextual factors.
The metadata parsing step extracts fields such as age/height/weight, gender, and yes/no questionnaire answers (coffee, sport, smoking, illness),
then encodes them into machine-readable columns (e.g., one-hot indicators).

\subsection{Why include metadata?}
Metadata can support:
\begin{itemize}
  \item \textbf{EDA interpretation}: identifying whether some differences are demographic rather than stress-related.
  \item \textbf{Confound analysis}: checking whether particular subject groups systematically differ in physiology.
  \item \textbf{Optional covariate modeling}: adding subject-level context features (with care to avoid evaluation leakage).
\end{itemize}

\section{Final dataset: \texttt{data/m14\_merged.csv}}
\subsection{What one row represents}
Each row corresponds to one \textbf{30-second window} from one subject, represented by engineered features plus a class label and metadata.

\subsection{Label meanings}
The labels used in \texttt{m14\_merged.csv} are:
\begin{center}
\begin{tabular}{cl}
\toprule
\textbf{Label value} & \textbf{Meaning} \\
\midrule
0 & amusement \\
1 & baseline \\
2 & stress \\
\bottomrule
\end{tabular}
\end{center}

\subsection{3--4 sample rows (subset of columns)}
The full dataset has many columns. Table~\ref{tab:samples} shows 4 windows using a representative subset of features.

\begin{table}[!ht]
  \centering
  \small
  \begin{tabular}{@{}rrrrrrrrr@{}}
    \toprule
    idx & subject & label & EDA\_mean & EDA\_phasic\_mean & BVP\_std & TEMP\_mean & Resp\_std & net\_acc\_mean \\
    \midrule
    0 & 2 & 1 & 1.398 & 1.829 & 107.679 & 35.817 & 0.011 & 0.026 \\
    1 & 2 & 1 & 1.210 & 2.111 & 118.776 & 35.798 & 0.021 & 0.029 \\
    2 & 2 & 1 & 1.011 & 0.153 & 42.202 & 35.713 & 0.011 & 0.033 \\
    3 & 2 & 1 & 0.775 & 0.178 & 41.619 & 35.701 & 0.003 & 0.018 \\
    \bottomrule
  \end{tabular}
  \caption{Sample windows from \texttt{m14\_merged.csv} (representative subset of columns).}
  \label{tab:samples}
\end{table}

\subsection{Column-by-column explanation}
Table~\ref{tab:columns} lists every column header and explains what it measures.

\renewcommand{\arraystretch}{1.1}
\begin{longtable}{@{}p{0.30\textwidth}p{0.65\textwidth}@{}}
\caption{\label{tab:columns}Column headers in \texttt{m14\_merged.csv} and their meaning.}\\
\toprule
\textbf{Column} & \textbf{Meaning} \\
\midrule
\endfirsthead
\toprule
\textbf{Column} & \textbf{Meaning} \\
\midrule
\endhead
\bottomrule
\endfoot

\texttt{(blank first column)} & CSV index column from export; not a feature. \\

\texttt{EDA\_mean} & Mean of wrist electrodermal activity (EDA) in the window. \\
\texttt{EDA\_std} & Standard deviation of EDA in the window. \\
\texttt{EDA\_min} & Minimum EDA in the window. \\
\texttt{EDA\_max} & Maximum EDA in the window. \\

\texttt{EDA\_phasic\_mean} & Mean of the phasic EDA component (fast responses) in the window. \\
\texttt{EDA\_phasic\_std} & Std of the phasic EDA component in the window. \\
\texttt{EDA\_phasic\_min} & Minimum of the phasic EDA component in the window. \\
\texttt{EDA\_phasic\_max} & Maximum of the phasic EDA component in the window. \\

\texttt{EDA\_smna\_mean} & Mean of the sparse driver component from EDA decomposition in the window. \\
\texttt{EDA\_smna\_std} & Std of the sparse driver component in the window. \\
\texttt{EDA\_smna\_min} & Minimum of the sparse driver component in the window. \\
\texttt{EDA\_smna\_max} & Maximum of the sparse driver component in the window. \\

\texttt{EDA\_tonic\_mean} & Mean of the tonic EDA component (slow drift / baseline level) in the window. \\
\texttt{EDA\_tonic\_std} & Std of the tonic EDA component in the window. \\
\texttt{EDA\_tonic\_min} & Minimum of the tonic EDA component in the window. \\
\texttt{EDA\_tonic\_max} & Maximum of the tonic EDA component in the window. \\

\texttt{BVP\_mean} & Mean of wrist blood volume pulse (BVP) in the window. \\
\texttt{BVP\_std} & Std of BVP in the window (often increases with motion artifacts). \\
\texttt{BVP\_min} & Minimum BVP in the window. \\
\texttt{BVP\_max} & Maximum BVP in the window. \\

\texttt{TEMP\_mean} & Mean wrist skin temperature in the window. \\
\texttt{TEMP\_std} & Std of temperature in the window. \\
\texttt{TEMP\_min} & Minimum temperature in the window. \\
\texttt{TEMP\_max} & Maximum temperature in the window. \\

\texttt{ACC\_x\_mean} & Mean wrist acceleration (x-axis) in the window. \\
\texttt{ACC\_x\_std} & Std wrist acceleration (x-axis) in the window. \\
\texttt{ACC\_x\_min} & Minimum wrist acceleration (x-axis) in the window. \\
\texttt{ACC\_x\_max} & Maximum wrist acceleration (x-axis) in the window. \\

\texttt{ACC\_y\_mean} & Mean wrist acceleration (y-axis) in the window. \\
\texttt{ACC\_y\_std} & Std wrist acceleration (y-axis) in the window. \\
\texttt{ACC\_y\_min} & Minimum wrist acceleration (y-axis) in the window. \\
\texttt{ACC\_y\_max} & Maximum wrist acceleration (y-axis) in the window. \\

\texttt{ACC\_z\_mean} & Mean wrist acceleration (z-axis) in the window. \\
\texttt{ACC\_z\_std} & Std wrist acceleration (z-axis) in the window. \\
\texttt{ACC\_z\_min} & Minimum wrist acceleration (z-axis) in the window. \\
\texttt{ACC\_z\_max} & Maximum wrist acceleration (z-axis) in the window. \\

\texttt{Resp\_mean} & Mean respiration signal in the window. \\
\texttt{Resp\_std} & Std respiration signal in the window. \\
\texttt{Resp\_min} & Minimum respiration signal in the window. \\
\texttt{Resp\_max} & Maximum respiration signal in the window. \\

\texttt{net\_acc\_mean} & Mean acceleration magnitude (derived from ACC x/y/z) in the window. \\
\texttt{net\_acc\_std} & Std acceleration magnitude in the window. \\
\texttt{net\_acc\_min} & Minimum acceleration magnitude in the window. \\
\texttt{net\_acc\_max} & Maximum acceleration magnitude in the window. \\

\texttt{BVP\_peak\_freq} & Dominant frequency of BVP in the window (periodogram peak). \\
\texttt{TEMP\_slope} & Linear temperature trend across the window. \\

\texttt{subject} & Subject identifier (integer). \\
\texttt{label} & Window label (0=amusement, 1=baseline, 2=stress). \\

\texttt{age} & Subject age. \\
\texttt{height} & Subject height. \\
\texttt{weight} & Subject weight. \\

\texttt{gender\_ female} & One-hot indicator: female. \\
\texttt{gender\_ male} & One-hot indicator: male. \\
\texttt{coffee\_today\_YES} & One-hot indicator: drank coffee today. \\
\texttt{sport\_today\_YES} & One-hot indicator: did sports today. \\
\texttt{smoker\_NO} & One-hot indicator: not a smoker. \\
\texttt{smoker\_YES} & One-hot indicator: smoker. \\
\texttt{feel\_ill\_today\_YES} & One-hot indicator: feels ill today. \\

\end{longtable}

\section{Exploratory data analysis (EDA) and why each plot is used}
The EDA is designed to answer three modeling-critical questions:
\begin{enumerate}
  \item \textbf{Is the dataset balanced and comparable across subjects?}
  \item \textbf{Do features separate the classes, and which modalities matter most?}
  \item \textbf{Are there artifacts/outliers or subject-specific effects that will hurt generalization?}
\end{enumerate}

\subsection{Dataset composition and class balance}
\begin{figure}[!ht]
  \centering
  \begin{subfigure}{0.49\linewidth}
    \centering
    \includegraphics[width=\linewidth]{\detokenize{eda_outputs_m14/01_label_counts.png}}
    \caption{Overall class balance.}
  \end{subfigure}
  \begin{subfigure}{0.49\linewidth}
    \centering
    \includegraphics[width=\linewidth]{\detokenize{eda_outputs_m14/02_windows_per_subject_stacked.png}}
    \caption{Per-subject window counts split by label.}
  \end{subfigure}
  \caption{Composition plots used to assess imbalance and protocol consistency.}
  \label{fig:composition}
\end{figure}

\paragraph{What and why.}
These plots quantify how many training examples exist per class and per subject.
They guide choices such as class-weighted loss, balanced sampling, and whether LOSO evaluation is necessary to prevent subject leakage.

\subsection{Feature separability by label (univariate comparisons)}
\begin{figure}[!ht]
  \centering
  \begin{subfigure}{0.49\linewidth}
    \centering
    \includegraphics[width=\linewidth]{\detokenize{eda_outputs_m14/03_box_eda_mean_by_label.png}}
    \caption{EDA mean by label.}
  \end{subfigure}
  \begin{subfigure}{0.49\linewidth}
    \centering
    \includegraphics[width=\linewidth]{\detokenize{eda_outputs_m14/04_box_eda_phasic_mean_by_label.png}}
    \caption{EDA phasic mean by label.}
  \end{subfigure}
  \caption{Univariate comparisons to see which physiological features shift across conditions.}
  \label{fig:univariate}
\end{figure}

\paragraph{What and why.}
Boxplots highlight whether a feature tends to be higher/lower under stress and how much overlap exists between classes.
This helps prioritize modalities for modeling (e.g., EDA-heavy vs multimodal fusion) and motivates feature scaling/robustification.

\subsection{Movement confounding check}
\begin{figure}[!ht]
  \centering
  \includegraphics[width=0.62\linewidth]{\detokenize{eda_outputs_m14/05_box_net_acc_mean_by_label.png}}
  \caption{Acceleration magnitude by label (used to detect motion confounding).}
  \label{fig:motion_conf}
\end{figure}

\paragraph{What and why.}
Wearable physiology can be corrupted by motion. If motion magnitude differs by label, a classifier may learn motion artifacts rather than stress physiology.
This motivates motion-aware preprocessing, removing low-quality windows, and ablation studies (physiology-only vs motion-only vs fused).

\subsection{Correlation and redundancy}
\begin{figure}[!ht]
  \centering
  \includegraphics[width=0.72\linewidth]{\detokenize{eda_outputs_m14/06_corr_heatmap_subset.png}}
  \caption{Correlation heatmap (subset) for redundancy/interaction analysis.}
  \label{fig:corr}
\end{figure}

\paragraph{What and why.}
Correlation reveals redundant features and interactions across modalities. This informs regularization strength, feature selection,
and whether dimensionality reduction (or learned representations) may help.

\subsection{Low-dimensional projections (PCA)}
\begin{figure}[!ht]
  \centering
  \begin{subfigure}{0.49\linewidth}
    \centering
    \includegraphics[width=\linewidth]{\detokenize{eda_outputs_m14/07_pca_scatter_by_label.png}}
    \caption{PCA colored by label.}
  \end{subfigure}
  \begin{subfigure}{0.49\linewidth}
    \centering
    \includegraphics[width=\linewidth]{\detokenize{eda_outputs_m14/08_pca_scatter_by_subject.png}}
    \caption{PCA colored by subject.}
  \end{subfigure}
  \caption{PCA projections to compare label structure vs subject structure.}
  \label{fig:pca}
\end{figure}

\paragraph{What and why.}
PCA provides a quick visual check of whether windows cluster more strongly by label or by subject.
Strong subject clustering is a common signal of cross-subject distribution shift and motivates LOSO evaluation and normalization strategies.

\subsection{Window timelines and outlier windows (within-subject)}
\begin{figure}[!ht]
  \centering
  \includegraphics[width=0.92\linewidth]{\detokenize{eda_outputs_m14_timeseries/subject_02_label_timeline.png}}
  \caption{Within-subject label segments across window index.}
  \label{fig:timeline}
\end{figure}

\begin{figure}[!ht]
  \centering
  \includegraphics[width=0.78\linewidth]{\detokenize{eda_outputs_m14_timeseries/subject_02_feature_traces.png}}
  \caption{Within-subject feature traces across windows with outliers highlighted.}
  \label{fig:traces}
\end{figure}

\paragraph{What and why.}
Treating windowed features as a sequence helps detect sudden jumps, drift, and rare outlier windows.
Outlier analysis informs robust scaling, clipping strategies, and data-cleaning rules prior to training.

\subsection{Raw signal inspection with label overlay}
\begin{figure}[!ht]
  \centering
  \includegraphics[width=0.92\linewidth]{\detokenize{eda_outputs_wesad_raw/S02_raw_EDA.png}}
  \caption{Raw EDA with protocol labels overlaid.}
  \label{fig:raw_eda}
\end{figure}

\begin{figure}[!ht]
  \centering
  \includegraphics[width=0.92\linewidth]{\detokenize{eda_outputs_wesad_raw/S02_raw_BVP.png}}
  \caption{Raw BVP with protocol labels overlaid (spikes often indicate motion artifacts).}
  \label{fig:raw_bvp}
\end{figure}

\begin{figure}[!ht]
  \centering
  \includegraphics[width=0.92\linewidth]{\detokenize{eda_outputs_wesad_raw/S02_raw_ACC_mag.png}}
  \caption{Raw acceleration magnitude with labels overlaid (used to validate motion artifacts).}
  \label{fig:raw_acc}
\end{figure}

\begin{figure}[!ht]
  \centering
  \includegraphics[width=0.92\linewidth]{\detokenize{eda_outputs_wesad_raw/S02_raw_TEMP.png}}
  \caption{Raw temperature with labels overlaid (slow trends and drift).}
  \label{fig:raw_temp}
\end{figure}

\begin{figure}[!ht]
  \centering
  \includegraphics[width=0.92\linewidth]{\detokenize{eda_outputs_wesad_raw/S02_raw_EDA_BVP_combined.png}}
  \caption{Combined raw EDA + BVP view with labels overlaid.}
  \label{fig:raw_combined}
\end{figure}

\paragraph{What and why.}
Window-level features can hide short-lived artifacts. Raw plots are used to visually confirm:
(i) whether BVP spikes align with movement bursts, (ii) whether some segments are consistently noisier than others, and
(iii) whether label timing matches the observed signal changes. This informs preprocessing and window-quality filtering.

\section{How the EDA informs model selection and training}
The transformation and EDA support the following decisions:
\begin{itemize}
  \item \textbf{Validation strategy}: prefer cross-subject evaluation (e.g., LOSO) to measure generalization.
  \item \textbf{Scaling and normalization}: use standardization (and consider per-subject normalization) due to subject-specific baselines.
  \item \textbf{Confound-aware modeling}: explicitly test motion-only vs physiology-only vs fused models to ensure the model learns stress physiology.
  \item \textbf{Robustness}: address outliers/artifacts through robust scaling, clipping, or removal of low-quality windows.
  \item \textbf{Feature/model choice}: if univariate plots show limited separation, use multimodal fusion and/or more expressive models.
\end{itemize}

\end{document}

\documentclass[11pt]{article}

\usepackage[margin=1in]{geometry}
\usepackage{graphicx}
\usepackage{subcaption}
\usepackage{booktabs}
\usepackage{longtable}
\usepackage{array}
\usepackage{xcolor}
\usepackage[hidelinks]{hyperref}

\title{WESAD: Data Transformation and Exploratory Data Analysis Report}
\author{Anishwar Chakraborty}
\date{\today}

\begin{document}
\maketitle

\section{What is WESAD and what does it contain?}
WESAD (Wearable Stress and Affect Detection) is a multimodal wearable-sensing dataset collected to study stress and affect.
Subjects follow a lab protocol that includes (at minimum) three key conditions:
\textbf{baseline} (rest), \textbf{stress} (stress induction), and \textbf{amusement} (positive affect).

\subsection{Modalities}
Each subject is stored as a Python pickle file (\texttt{.pkl}) containing:
\begin{itemize}
  \item \textbf{Wrist (Empatica E4)}: acceleration (ACC), blood volume pulse (BVP), electrodermal activity (EDA), and skin temperature (TEMP).
  \item \textbf{Chest (RespiBAN)}: signals such as respiration (Resp) and additional chest channels (e.g., ECG/EMG/EDA/Temp) depending on the recording.
  \item \textbf{A high-frequency label stream}: per-sample protocol labels indicating which condition the subject is in over time.
  \item \textbf{Subject-level metadata}: demographics and questionnaire-like fields contained in per-subject readme text files (e.g., age, gender, smoking).
\end{itemize}

\subsection{Why WESAD is suitable for modeling}
WESAD enables subject-independent stress detection because:
\begin{itemize}
  \item it provides \textbf{physiological} (EDA/BVP/TEMP/Resp) and \textbf{movement} (ACC) channels,
  \item it supports \textbf{window-based} feature extraction for classical ML or deep learning,
  \item it supports \textbf{leave-one-subject-out (LOSO)} evaluation for generalization across subjects.
\end{itemize}

\section{Data transformation: \texttt{data\_wrangling.py}}
The goal of \texttt{data\_wrangling.py} is to transform raw subject \texttt{.pkl} recordings into a single tabular dataset of
\textbf{fixed-length windows} with \textbf{engineered features} and \textbf{labels} suitable for EDA and model training.

\subsection{Inputs and sampling rates}
The script uses (see \texttt{fs\_dict} and \texttt{WINDOW\_IN\_SECONDS}):
\begin{itemize}
  \item Wrist ACC at 32 Hz, wrist BVP at 64 Hz, wrist EDA at 4 Hz, wrist TEMP at 4 Hz.
  \item A label stream at 700 Hz (and chest Resp is also treated as 700 Hz in this pipeline).
  \item Window length = 30 seconds.
\end{itemize}

\subsection{Step-by-step transformation}
For each subject \texttt{S\{id\}}:
\begin{enumerate}
  \item \textbf{Load raw data} via \texttt{SubjectData}, reading \texttt{data/WESAD/S\{id\}/S\{id\}.pkl}.
        The pipeline uses wrist sensors and additionally injects chest \texttt{Resp} into the wrist dictionary.
  \item \textbf{Preprocess and align signals} in \texttt{compute\_features}:
  \begin{itemize}
    \item EDA is low-pass filtered (to reduce high-frequency noise).
    \item ACC is filtered (FIR smoothing) to reduce jitter.
    \item Each sensor is assigned a time index based on its sampling rate, then all sensors are outer-joined into one frame.
    \item The high-frequency label stream is joined and then forward/backward filled to cover all timestamps.
    \item cvxEDA decomposition is applied to the filtered EDA to produce:
          \texttt{EDA\_phasic}, \texttt{EDA\_tonic}, and \texttt{EDA\_smna}.
  \end{itemize}
  \item \textbf{Split into protocol conditions} using the raw WESAD label convention inside \texttt{compute\_features}:
        baseline = 1, stress = 2, amusement = 3.
  \item \textbf{Window the aligned data} into 30s windows in \texttt{get\_samples}:
  \begin{itemize}
    \item Window length is computed in label-samples: $700 \times 30$.
    \item For each window, compute \texttt{net\_acc} (ACC magnitude) and then compute per-signal statistics:
          \texttt{mean}, \texttt{std}, \texttt{min}, \texttt{max}.
    \item Add two additional engineered features:
          \texttt{BVP\_peak\_freq} (dominant frequency) and \texttt{TEMP\_slope} (linear trend).
  \end{itemize}
  \item \textbf{Remap labels for modeling} in \texttt{make\_patient\_data}:
        windows are labeled as \texttt{baseline=1}, \texttt{stress=2}, \texttt{amusement=0}.
        (This is a project-specific remap; the raw label value for amusement is 3.)
  \item \textbf{Save per-subject feature CSVs} to \texttt{data/subject\_feats/S\{id\}\_feats\_4.csv}.
  \item \textbf{Combine subjects} in \texttt{combine\_files} into \texttt{data/may14\_feats4.csv}.
\end{enumerate}

\subsection{Why these transformations are useful for modeling}
\begin{itemize}
  \item \textbf{Windowing} converts variable-length signals into a supervised learning dataset (each row is a labeled example).
  \item \textbf{Summary features} make it easy to train baselines (MLP/trees/logistic regression) and perform fast EDA.
  \item \textbf{EDA phasic/tonic decomposition} helps separate fast arousal-related responses from slow drift.
  \item \textbf{Motion magnitude (net\_acc)} is critical to detect confounding from movement artifacts.
\end{itemize}

\section{Metadata parsing and final dataset: \texttt{readme\_parser.py}}
The goal of \texttt{readme\_parser.py} is to extract subject-level metadata from each \texttt{S\{id\}\_readme.txt},
one-hot encode categorical fields, and merge them into the window-feature dataset.

\subsection{What it extracts}
The parser searches for keys like \texttt{Age}, \texttt{Height}, \texttt{Weight}, \texttt{Gender}, and yes/no questionnaire fields
(coffee today, sport today, smoker, feel ill today). It then:
\begin{itemize}
  \item saves a combined metadata table (\texttt{data/WESAD/readmes.csv}),
  \item one-hot encodes categories using \texttt{pd.get\_dummies},
  \item constructs a numeric \texttt{subject} id from the folder name (e.g., ``S2'' $\rightarrow 2$),
  \item merges metadata with \texttt{data/may14\_feats4.csv} on \texttt{subject},
  \item outputs \texttt{data/m14\_merged.csv}.
\end{itemize}

\subsection{Why this matters}
Subject metadata can be used to:
\begin{itemize}
  \item perform stratified EDA (e.g., compare feature distributions by smoker status),
  \item analyze covariates and potential confounds,
  \item optionally augment models with subject-level covariates (with care to avoid leakage in evaluation).
\end{itemize}

\section{Final dataset: \texttt{data/m14\_merged.csv}}
\subsection{What one row represents}
Each row in \texttt{m14\_merged.csv} is one \textbf{30-second window} for a specific \textbf{subject}, with:
\begin{itemize}
  \item window-level sensor features (summary statistics and engineered features),
  \item a window label (target),
  \item subject-level metadata (repeated for all windows of that subject).
\end{itemize}

\subsection{Label meanings in \texttt{m14\_merged.csv}}
The project label mapping used in this CSV is:
\begin{center}
\begin{tabular}{cl}
\toprule
\textbf{Label value} & \textbf{Meaning} \\
\midrule
0 & amusement \\
1 & baseline \\
2 & stress \\
\bottomrule
\end{tabular}
\end{center}

\subsection{Sample rows (first 4 windows)}
For readability (the full CSV is very wide), Table~\ref{tab:samples} shows 4 sample windows with a small, representative subset of columns.

\begin{table}[!ht]
  \centering
  \small
  \begin{tabular}{@{}rrrrrrrrr@{}}
    \toprule
    idx & subject & label & EDA\_mean & EDA\_phasic\_mean & BVP\_std & TEMP\_mean & Resp\_std & net\_acc\_mean \\
    \midrule
    0 & 2 & 1 & 1.398 & 1.829 & 107.679 & 35.817 & 0.011 & 0.026 \\
    1 & 2 & 1 & 1.210 & 2.111 & 118.776 & 35.798 & 0.021 & 0.029 \\
    2 & 2 & 1 & 1.011 & 0.153 & 42.202 & 35.713 & 0.011 & 0.033 \\
    3 & 2 & 1 & 0.775 & 0.178 & 41.619 & 35.701 & 0.003 & 0.018 \\
    \bottomrule
  \end{tabular}
  \caption{Sample windows from \texttt{data/m14\_merged.csv} (subset of columns).}
  \label{tab:samples}
\end{table}

\subsection{Column-by-column description}
The columns can be grouped by (1) window-level sensor features, (2) engineered features, (3) identifiers/target, and (4) subject metadata.
Table~\ref{tab:columns} lists \emph{every} header and its meaning.

\renewcommand{\arraystretch}{1.1}
\begin{longtable}{@{}p{0.30\textwidth}p{0.65\textwidth}@{}}
\caption{\label{tab:columns}Column headers in \texttt{m14\_merged.csv} and their meaning.}\\
\toprule
\textbf{Column} & \textbf{Meaning} \\
\midrule
\endfirsthead
\toprule
\textbf{Column} & \textbf{Meaning} \\
\midrule
\endhead
\bottomrule
\endfoot

\texttt{(blank first column)} & CSV index saved by pandas when exporting; not a signal feature. \\

\texttt{EDA\_mean} & Mean of wrist EDA over the 30s window. \\
\texttt{EDA\_std} & Standard deviation of wrist EDA over the window. \\
\texttt{EDA\_min} & Minimum of wrist EDA over the window. \\
\texttt{EDA\_max} & Maximum of wrist EDA over the window. \\

\texttt{EDA\_phasic\_mean} & Mean of cvxEDA phasic component (fast SCR-related activity) over the window. \\
\texttt{EDA\_phasic\_std} & Std of cvxEDA phasic component over the window. \\
\texttt{EDA\_phasic\_min} & Min of cvxEDA phasic component over the window. \\
\texttt{EDA\_phasic\_max} & Max of cvxEDA phasic component over the window. \\

\texttt{EDA\_smna\_mean} & Mean of cvxEDA sparse driver (SMNA) component over the window. \\
\texttt{EDA\_smna\_std} & Std of cvxEDA SMNA component over the window. \\
\texttt{EDA\_smna\_min} & Min of cvxEDA SMNA component over the window. \\
\texttt{EDA\_smna\_max} & Max of cvxEDA SMNA component over the window. \\

\texttt{EDA\_tonic\_mean} & Mean of cvxEDA tonic component (slow baseline level) over the window. \\
\texttt{EDA\_tonic\_std} & Std of cvxEDA tonic component over the window. \\
\texttt{EDA\_tonic\_min} & Min of cvxEDA tonic component over the window. \\
\texttt{EDA\_tonic\_max} & Max of cvxEDA tonic component over the window. \\

\texttt{BVP\_mean} & Mean of wrist BVP over the window. \\
\texttt{BVP\_std} & Std of wrist BVP over the window (often increases with motion artifacts). \\
\texttt{BVP\_min} & Min of wrist BVP over the window. \\
\texttt{BVP\_max} & Max of wrist BVP over the window. \\

\texttt{TEMP\_mean} & Mean of wrist temperature over the window. \\
\texttt{TEMP\_std} & Std of wrist temperature over the window. \\
\texttt{TEMP\_min} & Min of wrist temperature over the window. \\
\texttt{TEMP\_max} & Max of wrist temperature over the window. \\

\texttt{ACC\_x\_mean} & Mean of wrist ACC x-axis over the window. \\
\texttt{ACC\_x\_std} & Std of wrist ACC x-axis over the window. \\
\texttt{ACC\_x\_min} & Min of wrist ACC x-axis over the window. \\
\texttt{ACC\_x\_max} & Max of wrist ACC x-axis over the window. \\

\texttt{ACC\_y\_mean} & Mean of wrist ACC y-axis over the window. \\
\texttt{ACC\_y\_std} & Std of wrist ACC y-axis over the window. \\
\texttt{ACC\_y\_min} & Min of wrist ACC y-axis over the window. \\
\texttt{ACC\_y\_max} & Max of wrist ACC y-axis over the window. \\

\texttt{ACC\_z\_mean} & Mean of wrist ACC z-axis over the window. \\
\texttt{ACC\_z\_std} & Std of wrist ACC z-axis over the window. \\
\texttt{ACC\_z\_min} & Min of wrist ACC z-axis over the window. \\
\texttt{ACC\_z\_max} & Max of wrist ACC z-axis over the window. \\

\texttt{Resp\_mean} & Mean of chest respiration signal (Resp) over the window. \\
\texttt{Resp\_std} & Std of chest respiration signal over the window. \\
\texttt{Resp\_min} & Min of chest respiration signal over the window. \\
\texttt{Resp\_max} & Max of chest respiration signal over the window. \\

\texttt{net\_acc\_mean} & Mean of acceleration magnitude (derived from ACC x/y/z) over the window. \\
\texttt{net\_acc\_std} & Std of acceleration magnitude over the window. \\
\texttt{net\_acc\_min} & Min of acceleration magnitude over the window. \\
\texttt{net\_acc\_max} & Max of acceleration magnitude over the window. \\

\texttt{BVP\_peak\_freq} & Dominant frequency of the BVP window (periodogram peak); proxy for periodicity/heart-rate band structure. \\
\texttt{TEMP\_slope} & Linear trend (slope) of temperature over the window. \\

\texttt{subject} & Subject id (integer). \\
\texttt{label} & Target class for the window (0=amusement, 1=baseline, 2=stress). \\

\texttt{age} & Subject age (from readme). \\
\texttt{height} & Subject height (from readme). \\
\texttt{weight} & Subject weight (from readme). \\

\texttt{gender\_ female} & One-hot: subject gender is female. \\
\texttt{gender\_ male} & One-hot: subject gender is male. \\
\texttt{coffee\_today\_YES} & One-hot: drank coffee today. \\
\texttt{sport\_today\_YES} & One-hot: did sports today. \\
\texttt{smoker\_NO} & One-hot: not a smoker. \\
\texttt{smoker\_YES} & One-hot: is a smoker. \\
\texttt{feel\_ill\_today\_YES} & One-hot: feels ill today. \\

\end{longtable}

\section{EDA scripts and plots (and how they inform modeling)}
This section explains each EDA script, the plots it generates, and how those plots are used for
model selection, training choices, and data-quality decisions.

\subsection{\texttt{eda\_m14.py}: dataset-level EDA on engineered windows}
\texttt{eda\_m14.py} loads \texttt{m14\_merged.csv}, maps numeric labels to names, selects sensor-window feature columns,
and produces dataset-level plots in \texttt{eda\_outputs\_m14/}.

\begin{figure}[!ht]
  \centering
  \begin{subfigure}{0.49\linewidth}
    \centering
    \includegraphics[width=\linewidth]{\detokenize{eda_outputs_m14/01_label_counts.png}}
    \caption{Class balance across all windows.}
  \end{subfigure}
  \begin{subfigure}{0.49\linewidth}
    \centering
    \includegraphics[width=\linewidth]{\detokenize{eda_outputs_m14/02_windows_per_subject_stacked.png}}
    \caption{Windows per subject, split by label.}
  \end{subfigure}
  \caption{Global dataset composition plots.}
  \label{fig:global_composition}
\end{figure}

\paragraph{Purpose and modeling use.}
\begin{itemize}
  \item Figure~\ref{fig:global_composition}(a) guides choices like class weighting, stratified sampling, and evaluation metrics.
  \item Figure~\ref{fig:global_composition}(b) confirms the protocol yields a comparable number of windows per subject and label,
        and motivates LOSO evaluation to avoid leakage across the same subject.
\end{itemize}

\begin{figure}[!ht]
  \centering
  \begin{subfigure}{0.49\linewidth}
    \centering
    \includegraphics[width=\linewidth]{\detokenize{eda_outputs_m14/03_box_eda_mean_by_label.png}}
    \caption{EDA mean by label.}
  \end{subfigure}
  \begin{subfigure}{0.49\linewidth}
    \centering
    \includegraphics[width=\linewidth]{\detokenize{eda_outputs_m14/04_box_eda_phasic_mean_by_label.png}}
    \caption{EDA phasic mean by label.}
  \end{subfigure}
  \caption{EDA-based separability checks.}
  \label{fig:eda_separability}
\end{figure}

\paragraph{Purpose and modeling use.}
\begin{itemize}
  \item If stress windows show systematically higher EDA (especially phasic activity), EDA is a high-value modality for the classifier.
  \item Large overlap suggests the need for multi-modal fusion (e.g., add BVP/Resp/TEMP/ACC) or more expressive models.
\end{itemize}

\begin{figure}[!ht]
  \centering
  \includegraphics[width=0.62\linewidth]{\detokenize{eda_outputs_m14/05_box_net_acc_mean_by_label.png}}
  \caption{Movement confound check: net acceleration magnitude by label.}
  \label{fig:net_acc}
\end{figure}

\paragraph{Purpose and modeling use.}
If net acceleration differs strongly by label, a model can ``cheat'' by learning movement rather than stress physiology.
This motivates: per-subject normalization, artifact handling, excluding extreme-motion windows, and/or running ablations
(physiology-only vs motion-only vs fused).

\begin{figure}[!ht]
  \centering
  \includegraphics[width=0.72\linewidth]{\detokenize{eda_outputs_m14/06_corr_heatmap_subset.png}}
  \caption{Feature correlation heatmap (subset).}
  \label{fig:corr}
\end{figure}

\paragraph{Purpose and modeling use.}
Correlation highlights redundancy and potential feature leakage:
highly correlated features may be removed (feature selection) or handled with regularization; it also informs dimensionality reduction.

\begin{figure}[!ht]
  \centering
  \begin{subfigure}{0.49\linewidth}
    \centering
    \includegraphics[width=\linewidth]{\detokenize{eda_outputs_m14/07_pca_scatter_by_label.png}}
    \caption{PCA colored by label.}
  \end{subfigure}
  \begin{subfigure}{0.49\linewidth}
    \centering
    \includegraphics[width=\linewidth]{\detokenize{eda_outputs_m14/08_pca_scatter_by_subject.png}}
    \caption{PCA colored by subject.}
  \end{subfigure}
  \caption{2D PCA projections of standardized window features.}
  \label{fig:pca}
\end{figure}

\paragraph{Purpose and modeling use.}
\begin{itemize}
  \item If samples separate mainly by \textbf{subject} rather than \textbf{label}, generalization is hard and LOSO is mandatory.
  \item Strong subject clustering motivates normalization strategies (global vs per-subject) and domain generalization techniques.
\end{itemize}

\subsection{\texttt{eda\_m14\_window\_timeseries.py}: window-index timelines and outlier windows}
\texttt{eda\_m14\_window\_timeseries.py} treats windows as a sequence (within a subject) and plots features vs window index.
It also exports top outlier windows based on a robust z-score (median/MAD).

\begin{figure}[!ht]
  \centering
  \includegraphics[width=0.92\linewidth]{\detokenize{eda_outputs_m14_timeseries/subject_02_label_timeline.png}}
  \caption{Subject 2: label segments across window index (order of windows for that subject).}
  \label{fig:label_timeline}
\end{figure}

\begin{figure}[!ht]
  \centering
  \includegraphics[width=0.78\linewidth]{\detokenize{eda_outputs_m14_timeseries/subject_02_feature_traces.png}}
  \caption{Subject 2: feature traces across windows with robust outliers marked (red points).}
  \label{fig:feature_traces}
\end{figure}

\paragraph{Purpose and modeling use.}
\begin{itemize}
  \item Identifies windows that are unusually high/low in key features (potential artifacts or rare events).
  \item Guides cleaning rules (drop extreme windows) or robust scaling (e.g., RobustScaler) before training.
  \item Helps verify label transitions and whether certain segments contain systematic feature shifts.
\end{itemize}

\subsection{\texttt{plot\_wesad\_raw\_timeseries.py}: raw signal inspection with label overlay}
\texttt{plot\_wesad\_raw\_timeseries.py} loads the raw \texttt{.pkl} recording and plots raw wrist signals
(EDA, BVP, TEMP, ACC magnitude) over real time, shading the background by protocol label.
This is used to visually diagnose spikes and artifacts that are not visible once data is summarized into windows.

\begin{figure}[!ht]
  \centering
  \includegraphics[width=0.92\linewidth]{\detokenize{eda_outputs_wesad_raw/S02_raw_EDA.png}}
  \caption{Raw wrist EDA for Subject 2 with baseline/stress/amusement overlay.}
  \label{fig:raw_eda}
\end{figure}

\begin{figure}[!ht]
  \centering
  \includegraphics[width=0.92\linewidth]{\detokenize{eda_outputs_wesad_raw/S02_raw_BVP.png}}
  \caption{Raw wrist BVP for Subject 2 with baseline/stress/amusement overlay (motion artifacts appear as large spikes).}
  \label{fig:raw_bvp}
\end{figure}

\begin{figure}[!ht]
  \centering
  \includegraphics[width=0.92\linewidth]{\detokenize{eda_outputs_wesad_raw/S02_raw_ACC_mag.png}}
  \caption{Raw wrist acceleration magnitude for Subject 2 with label overlay (used to confirm movement artifacts).}
  \label{fig:raw_acc}
\end{figure}

\begin{figure}[!ht]
  \centering
  \includegraphics[width=0.92\linewidth]{\detokenize{eda_outputs_wesad_raw/S02_raw_TEMP.png}}
  \caption{Raw wrist temperature for Subject 2 with label overlay (typically slow-varying; useful for drift).}
  \label{fig:raw_temp}
\end{figure}

\begin{figure}[!ht]
  \centering
  \includegraphics[width=0.92\linewidth]{\detokenize{eda_outputs_wesad_raw/S02_raw_EDA_BVP_combined.png}}
  \caption{Combined raw EDA + BVP view for Subject 2 with label overlay.}
  \label{fig:raw_combined}
\end{figure}

\paragraph{Purpose and modeling use.}
\begin{itemize}
  \item Confirms whether spikes in BVP co-occur with ACC bursts (a strong indicator of motion artifact).
  \item Motivates artifact-aware preprocessing (filtering, peak rejection) or dropping extreme windows.
  \item Ensures the label overlay aligns with protocol segments, preventing label/feature misalignment.
\end{itemize}

\section{How these results feed into model selection and training}
Based on the transformation and EDA outputs, the main modeling decisions supported by this report are:
\begin{itemize}
  \item \textbf{Evaluation}: use LOSO cross-validation because subject effects are strong (see PCA-by-subject).
  \item \textbf{Class imbalance}: consider class-weighted loss or balanced sampling if label counts differ.
  \item \textbf{Feature scaling}: use standardization (and potentially per-subject normalization) due to cross-subject distribution shift.
  \item \textbf{Confound control}: explicitly analyze motion (net\_acc, ACC) and run modality ablations to ensure the model learns stress physiology.
  \item \textbf{Robustness}: handle outliers via robust scaling, clipping, or window-quality filtering informed by raw-signal inspection.
\end{itemize}

\end{document}

\\documentclass[11pt]{article}

\\usepackage[margin=1in]{geometry}
\\usepackage{graphicx}
\\usepackage{subcaption}
\\usepackage{booktabs}
\\usepackage{longtable}
\\usepackage{array}
\\usepackage{xcolor}
\\usepackage[hidelinks]{hyperref}

\\title{WESAD: Data Transformation and Exploratory Data Analysis Report}
\\author{Anishwar Chakraborty}
\\date{\\today}

\\begin{document}
\\maketitle

\\section{What is WESAD and what does it contain?}
WESAD (Wearable Stress and Affect Detection) is a multimodal wearable-sensing dataset collected to study stress and affect.
Subjects follow a lab protocol that includes (at minimum) three key conditions:
\\textbf{baseline} (rest), \\textbf{stress} (stress induction), and \\textbf{amusement} (positive affect).

\\subsection{Modalities}
Each subject is stored as a Python pickle file (\\texttt{.pkl}) containing:
\\begin{itemize}
  \\item \\textbf{Wrist (Empatica E4)}: acceleration (ACC), blood volume pulse (BVP), electrodermal activity (EDA), and skin temperature (TEMP).
  \\item \\textbf{Chest (RespiBAN)}: signals such as respiration (Resp) and additional chest channels (e.g., ECG/EMG/EDA/Temp) depending on the recording.
  \\item \\textbf{A high-frequency label stream}: per-sample protocol labels indicating which condition the subject is in over time.
  \\item \\textbf{Subject-level metadata}: demographics and questionnaire-like fields contained in per-subject readme text files (e.g., age, gender, smoking).
\\end{itemize}

\\subsection{Why WESAD is suitable for modeling}
WESAD enables subject-independent stress detection because:
\\begin{itemize}
  \\item it provides \\textbf{physiological} (EDA/BVP/TEMP/Resp) and \\textbf{movement} (ACC) channels,
  \\item it supports \\textbf{window-based} feature extraction for classical ML or deep learning,
  \\item it supports \\textbf{leave-one-subject-out (LOSO)} evaluation for generalization across subjects.
\\end{itemize}

\\section{Data transformation: \\texttt{data\\_wrangling.py}}
The goal of \\texttt{data\\_wrangling.py} is to transform raw subject \\texttt{.pkl} recordings into a single tabular dataset of
\\textbf{fixed-length windows} with \\textbf{engineered features} and \\textbf{labels} suitable for EDA and model training.

\\subsection{Inputs and sampling rates}
The script uses (see \\texttt{fs\\_dict} and \\texttt{WINDOW\\_IN\\_SECONDS}):
\\begin{itemize}
  \\item Wrist ACC at 32 Hz, wrist BVP at 64 Hz, wrist EDA at 4 Hz, wrist TEMP at 4 Hz.
  \\item A label stream at 700 Hz (and chest Resp is also treated as 700 Hz in this pipeline).
  \\item Window length = 30 seconds.
\\end{itemize}

\\subsection{Step-by-step transformation}
For each subject \\texttt{S\\{id\\}}:
\\begin{enumerate}
  \\item \\textbf{Load raw data} via \\texttt{SubjectData}, reading \\texttt{data/WESAD/S\\{id\\}/S\\{id\\}.pkl}.
        The pipeline uses wrist sensors and additionally injects chest \\texttt{Resp} into the wrist dictionary.
  \\item \\textbf{Preprocess and align signals} in \\texttt{compute\\_features}:
  \\begin{itemize}
    \\item EDA is low-pass filtered (to reduce high-frequency noise).
    \\item ACC is filtered (FIR smoothing) to reduce jitter.
    \\item Each sensor is assigned a time index based on its sampling rate, then all sensors are outer-joined into one frame.
    \\item The high-frequency label stream is joined and then forward/backward filled to cover all timestamps.
    \\item cvxEDA decomposition is applied to the filtered EDA to produce:
          \\texttt{EDA\\_phasic}, \\texttt{EDA\\_tonic}, and \\texttt{EDA\\_smna}.
  \\end{itemize}
  \\item \\textbf{Split into protocol conditions} using the raw WESAD label convention inside \\texttt{compute\\_features}:
        baseline = 1, stress = 2, amusement = 3.
  \\item \\textbf{Window the aligned data} into 30s windows in \\texttt{get\\_samples}:
  \\begin{itemize}
    \\item Window length is computed in label-samples: \\(700 \\times 30\\).
    \\item For each window, compute \\texttt{net\\_acc} (ACC magnitude) and then compute per-signal statistics:
          \\texttt{mean}, \\texttt{std}, \\texttt{min}, \\texttt{max}.
    \\item Add two additional engineered features:
          \\texttt{BVP\\_peak\\_freq} (dominant frequency) and \\texttt{TEMP\\_slope} (linear trend).
  \\end{itemize}
  \\item \\textbf{Remap labels for modeling} in \\texttt{make\\_patient\\_data}:
        windows are labeled as \\texttt{baseline=1}, \\texttt{stress=2}, \\texttt{amusement=0}.
        (This is a project-specific remap; the raw label value for amusement is 3.)
  \\item \\textbf{Save per-subject feature CSVs} to \\texttt{data/subject\\_feats/S\\{id\\}\\_feats\\_4.csv}.
  \\item \\textbf{Combine subjects} in \\texttt{combine\\_files} into \\texttt{data/may14\\_feats4.csv}.
\\end{enumerate}

\\subsection{Why these transformations are useful for modeling}
\\begin{itemize}
  \\item \\textbf{Windowing} converts variable-length signals into a supervised learning dataset (each row is a labeled example).
  \\item \\textbf{Summary features} make it easy to train baselines (MLP/trees/logistic regression) and perform fast EDA.
  \\item \\textbf{EDA phasic/tonic decomposition} helps separate fast arousal-related responses from slow drift.
  \\item \\textbf{Motion magnitude (net\\_acc)} is critical to detect confounding from movement artifacts.
\\end{itemize}

\\section{Metadata parsing and final dataset: \\texttt{readme\\_parser.py}}
The goal of \\texttt{readme\\_parser.py} is to extract subject-level metadata from each \\texttt{S\\{id\\}\\_readme.txt},
one-hot encode categorical fields, and merge them into the window-feature dataset.

\\subsection{What it extracts}
The parser searches for keys like \\texttt{Age}, \\texttt{Height}, \\texttt{Weight}, \\texttt{Gender}, and yes/no questionnaire fields
(coffee today, sport today, smoker, feel ill today). It then:
\\begin{itemize}
  \\item saves a combined metadata table (\\texttt{data/WESAD/readmes.csv}),
  \\item one-hot encodes categories using \\texttt{pd.get\\_dummies},
  \\item constructs a numeric \\texttt{subject} id from the folder name (e.g., ``S2'' \\(\\rightarrow 2\\)),
  \\item merges metadata with \\texttt{data/may14\\_feats4.csv} on \\texttt{subject},
  \\item outputs \\texttt{data/m14\\_merged.csv}.
\\end{itemize}

\\subsection{Why this matters}
Subject metadata can be used to:
\\begin{itemize}
  \\item perform stratified EDA (e.g., compare feature distributions by smoker status),
  \\item analyze covariates and potential confounds,
  \\item optionally augment models with subject-level covariates (with care to avoid leakage in evaluation).
\\end{itemize}

\\section{Final dataset: \\texttt{data/m14\\_merged.csv}}
\\subsection{What one row represents}
Each row in \\texttt{m14\\_merged.csv} is one \\textbf{30-second window} for a specific \\textbf{subject}, with:
\\begin{itemize}
  \\item window-level sensor features (summary statistics and engineered features),
  \\item a window label (target),
  \\item subject-level metadata (repeated for all windows of that subject).
\\end{itemize}

\\subsection{Label meanings in \\texttt{m14\\_merged.csv}}
The project label mapping used in this CSV is:
\\begin{center}
\\begin{tabular}{cl}
\\toprule
\\textbf{Label value} & \\textbf{Meaning} \\\\
\\midrule
0 & amusement \\\\
1 & baseline \\\\
2 & stress \\\\
\\bottomrule
\\end{tabular}
\\end{center}

\\subsection{Sample rows (first 4 windows)}
The following are sample lines (header + 4 rows) directly from \\texttt{data/m14\\_merged.csv}.
They are shown verbatim to preserve the wide set of columns.
\\begin{small}
\\begin{verbatim}
,EDA_mean,EDA_std,EDA_min,EDA_max,EDA_phasic_mean,EDA_phasic_std,EDA_phasic_min,EDA_phasic_max,EDA_smna_mean,EDA_smna_std,EDA_smna_min,EDA_smna_max,EDA_tonic_mean,EDA_tonic_std,EDA_tonic_min,EDA_tonic_max,BVP_mean,BVP_std,BVP_min,BVP_max,TEMP_mean,TEMP_std,TEMP_min,TEMP_max,ACC_x_mean,ACC_x_std,ACC_x_min,ACC_x_max,ACC_y_mean,ACC_y_std,ACC_y_min,ACC_y_max,ACC_z_mean,ACC_z_std,ACC_z_min,ACC_z_max,Resp_mean,Resp_std,Resp_min,Resp_max,net_acc_mean,net_acc_std,net_acc_min,net_acc_max,BVP_peak_freq,TEMP_slope,subject,label,age,height,weight,gender_ female,gender_ male,coffee_today_YES,sport_today_YES,smoker_NO,smoker_YES,feel_ill_today_YES
0,1.3979684884184072,0.1421283950003331,1.109298662480257,1.678399464596154,1.8286466800415988,1.0941151699995388,0.3703718448675747,4.327158402584045,1.2872297863894322,1.96107552424159,5.86535549018663e-08,11.710735943358904,1.2278044130911083,1.0017936627320914,-0.5999062719487944,2.550225067587716,-0.181673261729791,107.67879846954864,-358.13,554.77,35.81709090909091,0.0127319186087378,35.79,35.84,2.366352592827641e-05,8.94181485553418e-06,-1.9310795279339348e-05,7.007060001360278e-05,0.0214135759067394,0.0170525491050647,-0.0412430914886915,0.0522686511839442,2.030445688277317e-05,7.672512318672817e-06,-1.656960215099968e-05,6.012398494791312e-05,0.0295100101555683,0.011151045199089,-0.024081861956356,0.0873827562416348,0.0256889686507238,0.0094435480815385,5.33266539833442e-05,0.0522686573775084,0.1356698699830412,-0.0001690598020873,2,1,27,175,80,False,True,False,False,True,False,False
1,1.2101317674035457,0.092299131682188,1.014137872636249,1.4857995814446945,2.1109571853194464,1.2291722511400092,0.540783574244884,4.461295191018592,1.4692771168315963,2.867105708678988,3.4907387203020335e-08,17.430256888008586,0.3758003884346144,1.1776194651315297,-1.2149235778082526,1.8705673636050693,-0.8301470588235295,118.77568459744305,-392.28,438.16,35.79756756756757,0.0300368324263933,35.75,35.87,1.3913883618691409e-05,1.6702213411614888e-05,-3.034553543896184e-05,5.738064883003692e-05,0.0290941071450154,0.0159089807077325,-0.0047809503390381,0.0529974731597957,1.1938789293866704e-05,1.4331311956252988e-05,-2.603794623728521e-05,4.923538924868476e-05,0.017351549728282,0.020828784725114,-0.0378429259314166,0.0715575326703151,0.0291658493082437,0.0157769470810625,5.431418461266541e-05,0.0529974765306299,0.0950226244343891,-0.0007886977886977,2,1,27,175,80,False,True,False,False,True,False,False
2,1.010977164701998,0.1027831346412562,0.8322161167473113,1.1909666436032378,0.1530671558713177,0.12971226066199,0.0069728664445025,0.5452556472237626,0.1052543513378176,0.2464066245504445,5.303034544823505e-08,1.301887673913666,1.7274534403513746,0.2947030266481896,1.1367540494971722,2.0369006233903604,0.9396834369700398,42.20196927852661,-240.61,209.89,35.71290909090909,0.0277388031645353,35.66,35.75,1.6710461877882547e-05,8.853249473518395e-06,2.206948031924497e-06,4.358722363050881e-05,0.0322421231909387,0.0139468930490651,-0.0102793608635888,0.04198394359549,1.4338389541739969e-05,7.596518910678744e-06,1.8936688172571063e-06,3.739995914082785e-05,0.0208390711179392,0.0110405982043918,0.0027522127950121,0.0543562027014893,0.0326904451685088,0.0128597363533526,5.5301715241986593e-05,0.0419839458678118,0.07687959299039,-0.0007169938913058,2,1,27,175,80,False,True,False,False,True,False,False
3,0.7751874460203967,0.0466014846677679,0.6939961352625406,0.876819417928938,0.1782723846197622,0.1271378479148265,0.0028143930883128,0.362001176673944,0.1112683500138614,0.201946303798118,3.1350842810695674e-08,1.120160175330418,0.9872456458692074,0.0426408310487595,0.9123845529294374,1.127028011983083,0.1074038461538463,41.61864338063139,-289.26,145.36,35.70081081081081,0.0195929085100219,35.66,35.73,2.762367447538074e-05,2.5555057146194463e-06,1.1034740159622484e-05,3.255248347088633e-05,0.0182201709604528,0.0041485833031746,0.0041637412447512,0.0330512821426477,2.370245706532304e-05,2.1927482723176064e-06,9.468344086285533e-06,2.793161505454232e-05,0.0344485820402832,0.0031868876945729,0.0137610639750605,0.0405951387264287,0.0182202107540083,0.0041485695609252,0.0041639195255542,0.0330512930887862,0.1402714932126696,7.528957528957787e-05,2,1,27,175,80,False,True,False,False,True,False,False
\\end{verbatim}
\\end{small}

\\subsection{Column-by-column description}
The columns can be grouped by (1) window-level sensor features, (2) engineered features, (3) identifiers/target, and (4) subject metadata.
Table~\\ref{tab:columns} lists \\emph{every} header and its meaning.

\\renewcommand{\\arraystretch}{1.1}
\\begin{longtable}{@{}p{0.29\\textwidth}p{0.66\\textwidth}@{}}
\\caption{\\label{tab:columns}Column headers in \\texttt{m14\\_merged.csv} and their meaning.}\\\\
\\toprule
\\textbf{Column} & \\textbf{Meaning} \\\\
\\midrule
\\endfirsthead
\\toprule
\\textbf{Column} & \\textbf{Meaning} \\\\
\\midrule
\\endhead
\\bottomrule
\\endfoot

\\texttt{(blank first column)} & CSV index saved by pandas when exporting; not a signal feature. \\\\

\\texttt{EDA\\_mean} & Mean of wrist EDA over the 30s window. \\\\
\\texttt{EDA\\_std} & Standard deviation of wrist EDA over the window. \\\\
\\texttt{EDA\\_min} & Minimum of wrist EDA over the window. \\\\
\\texttt{EDA\\_max} & Maximum of wrist EDA over the window. \\\\

\\texttt{EDA\\_phasic\\_mean} & Mean of cvxEDA phasic component (fast SCR-related activity) over the window. \\\\
\\texttt{EDA\\_phasic\\_std} & Std of cvxEDA phasic component over the window. \\\\
\\texttt{EDA\\_phasic\\_min} & Min of cvxEDA phasic component over the window. \\\\
\\texttt{EDA\\_phasic\\_max} & Max of cvxEDA phasic component over the window. \\\\

\\texttt{EDA\\_smna\\_mean} & Mean of cvxEDA sparse driver (SMNA) component over the window. \\\\
\\texttt{EDA\\_smna\\_std} & Std of cvxEDA SMNA component over the window. \\\\
\\texttt{EDA\\_smna\\_min} & Min of cvxEDA SMNA component over the window. \\\\
\\texttt{EDA\\_smna\\_max} & Max of cvxEDA SMNA component over the window. \\\\

\\texttt{EDA\\_tonic\\_mean} & Mean of cvxEDA tonic component (slow baseline level) over the window. \\\\
\\texttt{EDA\\_tonic\\_std} & Std of cvxEDA tonic component over the window. \\\\
\\texttt{EDA\\_tonic\\_min} & Min of cvxEDA tonic component over the window. \\\\
\\texttt{EDA\\_tonic\\_max} & Max of cvxEDA tonic component over the window. \\\\

\\texttt{BVP\\_mean} & Mean of wrist BVP over the window. \\\\
\\texttt{BVP\\_std} & Std of wrist BVP over the window (often increases with motion artifacts). \\\\
\\texttt{BVP\\_min} & Min of wrist BVP over the window. \\\\
\\texttt{BVP\\_max} & Max of wrist BVP over the window. \\\\

\\texttt{TEMP\\_mean} & Mean of wrist temperature over the window. \\\\
\\texttt{TEMP\\_std} & Std of wrist temperature over the window. \\\\
\\texttt{TEMP\\_min} & Min of wrist temperature over the window. \\\\
\\texttt{TEMP\\_max} & Max of wrist temperature over the window. \\\\

\\texttt{ACC\\_x\\_mean} & Mean of wrist ACC x-axis over the window. \\\\
\\texttt{ACC\\_x\\_std} & Std of wrist ACC x-axis over the window. \\\\
\\texttt{ACC\\_x\\_min} & Min of wrist ACC x-axis over the window. \\\\
\\texttt{ACC\\_x\\_max} & Max of wrist ACC x-axis over the window. \\\\

\\texttt{ACC\\_y\\_mean} & Mean of wrist ACC y-axis over the window. \\\\
\\texttt{ACC\\_y\\_std} & Std of wrist ACC y-axis over the window. \\\\
\\texttt{ACC\\_y\\_min} & Min of wrist ACC y-axis over the window. \\\\
\\texttt{ACC\\_y\\_max} & Max of wrist ACC y-axis over the window. \\\\

\\texttt{ACC\\_z\\_mean} & Mean of wrist ACC z-axis over the window. \\\\
\\texttt{ACC\\_z\\_std} & Std of wrist ACC z-axis over the window. \\\\
\\texttt{ACC\\_z\\_min} & Min of wrist ACC z-axis over the window. \\\\
\\texttt{ACC\\_z\\_max} & Max of wrist ACC z-axis over the window. \\\\

\\texttt{Resp\\_mean} & Mean of chest respiration signal (Resp) over the window. \\\\
\\texttt{Resp\\_std} & Std of chest respiration signal over the window. \\\\
\\texttt{Resp\\_min} & Min of chest respiration signal over the window. \\\\
\\texttt{Resp\\_max} & Max of chest respiration signal over the window. \\\\

\\texttt{net\\_acc\\_mean} & Mean of acceleration magnitude (derived from ACC x/y/z) over the window. \\\\
\\texttt{net\\_acc\\_std} & Std of acceleration magnitude over the window. \\\\
\\texttt{net\\_acc\\_min} & Min of acceleration magnitude over the window. \\\\
\\texttt{net\\_acc\\_max} & Max of acceleration magnitude over the window. \\\\

\\texttt{BVP\\_peak\\_freq} & Dominant frequency of the BVP window (periodogram peak); proxy for periodicity/heart-rate band structure. \\\\
\\texttt{TEMP\\_slope} & Linear trend (slope) of temperature over the window. \\\\

\\texttt{subject} & Subject id (integer). \\\\
\\texttt{label} & Target class for the window (0=amusement, 1=baseline, 2=stress). \\\\

\\texttt{age} & Subject age (from readme). \\\\
\\texttt{height} & Subject height (from readme). \\\\
\\texttt{weight} & Subject weight (from readme). \\\\

\\texttt{gender\\_ female} & One-hot: subject gender is female. \\\\
\\texttt{gender\\_ male} & One-hot: subject gender is male. \\\\
\\texttt{coffee\\_today\\_YES} & One-hot: drank coffee today. \\\\
\\texttt{sport\\_today\\_YES} & One-hot: did sports today. \\\\
\\texttt{smoker\\_NO} & One-hot: not a smoker. \\\\
\\texttt{smoker\\_YES} & One-hot: is a smoker. \\\\
\\texttt{feel\\_ill\\_today\\_YES} & One-hot: feels ill today. \\\\

\\end{longtable}

\\section{EDA scripts and plots (and how they inform modeling)}
This section explains each EDA script, the plots it generates, and how those plots are used for
model selection, training choices, and data-quality decisions.

\\subsection{\\texttt{eda\\_m14.py}: dataset-level EDA on engineered windows}
\\texttt{eda\\_m14.py} loads \\texttt{m14\\_merged.csv}, maps numeric labels to names, selects sensor-window feature columns,
and produces dataset-level plots in \\texttt{eda\\_outputs\\_m14/}.

\\begin{figure}[!ht]
  \\centering
  \\begin{subfigure}{0.49\\linewidth}
    \\centering
    \\includegraphics[width=\\linewidth]{\\detokenize{eda_outputs_m14/01_label_counts.png}}
    \\caption{Class balance across all windows.}
  \\end{subfigure}
  \\begin{subfigure}{0.49\\linewidth}
    \\centering
    \\includegraphics[width=\\linewidth]{\\detokenize{eda_outputs_m14/02_windows_per_subject_stacked.png}}
    \\caption{Windows per subject, split by label.}
  \\end{subfigure}
  \\caption{Global dataset composition plots.}
  \\label{fig:global_composition}
\\end{figure}

\\paragraph{Purpose and modeling use.}
\\begin{itemize}
  \\item Figure~\\ref{fig:global_composition}(a) guides choices like class weighting, stratified sampling, and evaluation metrics.
  \\item Figure~\\ref{fig:global_composition}(b) confirms the protocol yields a comparable number of windows per subject and label,
        and motivates LOSO evaluation to avoid leakage across the same subject.
\\end{itemize}

\\begin{figure}[!ht]
  \\centering
  \\begin{subfigure}{0.49\\linewidth}
    \\centering
    \\includegraphics[width=\\linewidth]{\\detokenize{eda_outputs_m14/03_box_eda_mean_by_label.png}}
    \\caption{EDA mean by label.}
  \\end{subfigure}
  \\begin{subfigure}{0.49\\linewidth}
    \\centering
    \\includegraphics[width=\\linewidth]{\\detokenize{eda_outputs_m14/04_box_eda_phasic_mean_by_label.png}}
    \\caption{EDA phasic mean by label.}
  \\end{subfigure}
  \\caption{EDA-based separability checks.}
  \\label{fig:eda_separability}
\\end{figure}

\\paragraph{Purpose and modeling use.}
\\begin{itemize}
  \\item If stress windows show systematically higher EDA (especially phasic activity), EDA is a high-value modality for the classifier.
  \\item Large overlap suggests the need for multi-modal fusion (e.g., add BVP/Resp/TEMP/ACC) or more expressive models.
\\end{itemize}

\\begin{figure}[!ht]
  \\centering
  \\includegraphics[width=0.62\\linewidth]{\\detokenize{eda_outputs_m14/05_box_net_acc_mean_by_label.png}}
  \\caption{Movement confound check: net acceleration magnitude by label.}
  \\label{fig:net_acc}
\\end{figure}

\\paragraph{Purpose and modeling use.}
If net acceleration differs strongly by label, a model can ``cheat'' by learning movement rather than stress physiology.
This motivates: per-subject normalization, artifact handling, excluding extreme-motion windows, and/or running ablations
(physiology-only vs motion-only vs fused).

\\begin{figure}[!ht]
  \\centering
  \\includegraphics[width=0.72\\linewidth]{\\detokenize{eda_outputs_m14/06_corr_heatmap_subset.png}}
  \\caption{Feature correlation heatmap (subset).}
  \\label{fig:corr}
\\end{figure}

\\paragraph{Purpose and modeling use.}
Correlation highlights redundancy and potential feature leakage:
highly correlated features may be removed (feature selection) or handled with regularization; it also informs dimensionality reduction.

\\begin{figure}[!ht]
  \\centering
  \\begin{subfigure}{0.49\\linewidth}
    \\centering
    \\includegraphics[width=\\linewidth]{\\detokenize{eda_outputs_m14/07_pca_scatter_by_label.png}}
    \\caption{PCA colored by label.}
  \\end{subfigure}
  \\begin{subfigure}{0.49\\linewidth}
    \\centering
    \\includegraphics[width=\\linewidth]{\\detokenize{eda_outputs_m14/08_pca_scatter_by_subject.png}}
    \\caption{PCA colored by subject.}
  \\end{subfigure}
  \\caption{2D PCA projections of standardized window features.}
  \\label{fig:pca}
\\end{figure}

\\paragraph{Purpose and modeling use.}
\\begin{itemize}
  \\item If samples separate mainly by \\textbf{subject} rather than \\textbf{label}, generalization is hard and LOSO is mandatory.
  \\item Strong subject clustering motivates normalization strategies (global vs per-subject) and domain generalization techniques.
\\end{itemize}

\\subsection{\\texttt{eda\\_m14\\_window\\_timeseries.py}: window-index timelines and outlier windows}
\\texttt{eda\\_m14\\_window\\_timeseries.py} treats windows as a sequence (within a subject) and plots features vs window index.
It also exports top outlier windows based on a robust z-score (median/MAD).

\\begin{figure}[!ht]
  \\centering
  \\includegraphics[width=0.92\\linewidth]{\\detokenize{eda_outputs_m14_timeseries/subject_02_label_timeline.png}}
  \\caption{Subject 2: label segments across window index (order of windows for that subject).}
  \\label{fig:label_timeline}
\\end{figure}

\\begin{figure}[!ht]
  \\centering
  \\includegraphics[width=0.78\\linewidth]{\\detokenize{eda_outputs_m14_timeseries/subject_02_feature_traces.png}}
  \\caption{Subject 2: feature traces across windows with robust outliers marked (red points).}
  \\label{fig:feature_traces}
\\end{figure}

\\paragraph{Purpose and modeling use.}
\\begin{itemize}
  \\item Identifies windows that are unusually high/low in key features (potential artifacts or rare events).
  \\item Guides cleaning rules (drop extreme windows) or robust scaling (e.g., RobustScaler) before training.
  \\item Helps verify label transitions and whether certain segments contain systematic feature shifts.
\\end{itemize}

\\subsection{\\texttt{plot\\_wesad\\_raw\\_timeseries.py}: raw signal inspection with label overlay}
\\texttt{plot\\_wesad\\_raw\\_timeseries.py} loads the raw \\texttt{.pkl} recording and plots raw wrist signals
(EDA, BVP, TEMP, ACC magnitude) over real time, shading the background by protocol label.
This is used to visually diagnose spikes and artifacts that are not visible once data is summarized into windows.

\\begin{figure}[!ht]
  \\centering
  \\includegraphics[width=0.92\\linewidth]{\\detokenize{eda_outputs_wesad_raw/S02_raw_EDA.png}}
  \\caption{Raw wrist EDA for Subject 2 with baseline/stress/amusement overlay.}
  \\label{fig:raw_eda}
\\end{figure}

\\begin{figure}[!ht]
  \\centering
  \\includegraphics[width=0.92\\linewidth]{\\detokenize{eda_outputs_wesad_raw/S02_raw_BVP.png}}
  \\caption{Raw wrist BVP for Subject 2 with baseline/stress/amusement overlay (motion artifacts appear as large spikes).}
  \\label{fig:raw_bvp}
\\end{figure}

\\begin{figure}[!ht]
  \\centering
  \\includegraphics[width=0.92\\linewidth]{\\detokenize{eda_outputs_wesad_raw/S02_raw_ACC_mag.png}}
  \\caption{Raw wrist acceleration magnitude for Subject 2 with label overlay (used to confirm movement artifacts).}
  \\label{fig:raw_acc}
\\end{figure}

\\begin{figure}[!ht]
  \\centering
  \\includegraphics[width=0.92\\linewidth]{\\detokenize{eda_outputs_wesad_raw/S02_raw_TEMP.png}}
  \\caption{Raw wrist temperature for Subject 2 with label overlay (typically slow-varying; useful for drift).}
  \\label{fig:raw_temp}
\\end{figure}

\\begin{figure}[!ht]
  \\centering
  \\includegraphics[width=0.92\\linewidth]{\\detokenize{eda_outputs_wesad_raw/S02_raw_EDA_BVP_combined.png}}
  \\caption{Combined raw EDA + BVP view for Subject 2 with label overlay.}
  \\label{fig:raw_combined}
\\end{figure}

\\paragraph{Purpose and modeling use.}
\\begin{itemize}
  \\item Confirms whether spikes in BVP co-occur with ACC bursts (a strong indicator of motion artifact).
  \\item Motivates artifact-aware preprocessing (filtering, peak rejection) or dropping extreme windows.
  \\item Ensures the label overlay aligns with protocol segments, preventing label/feature misalignment.
\\end{itemize}

\\section{How these results feed into model selection and training}
Based on the transformation and EDA outputs, the main modeling decisions supported by this report are:
\\begin{itemize}
  \\item \\textbf{Evaluation}: use LOSO cross-validation because subject effects are strong (see PCA-by-subject).
  \\item \\textbf{Class imbalance}: consider class-weighted loss or balanced sampling if label counts differ.
  \\item \\textbf{Feature scaling}: use standardization (and potentially per-subject normalization) due to cross-subject distribution shift.
  \\item \\textbf{Confound control}: explicitly analyze motion (net\\_acc, ACC) and run modality ablations to ensure the model learns stress physiology.
  \\item \\textbf{Robustness}: handle outliers via robust scaling, clipping, or window-quality filtering informed by raw-signal inspection.
\\end{itemize}

\\end{document}

